\documentclass{article}
\usepackage[fontsize=10bp]{fontsize}
\usepackage[utf8]{inputenc}
\usepackage[margin=1.0in]{geometry}
\usepackage{amsmath}
\usepackage{amssymb}
\usepackage{amsthm}
\usepackage{subfigure}
\usepackage[parfill]{parskip}
\usepackage{graphicx}
\usepackage[usenames,dvipsnames]{xcolor}
\usepackage{float}
\usepackage{amsfonts}
\usepackage{tcolorbox}
\usepackage{multicol}
\usepackage{wrapfig}
\usepackage{csquotes}
\usepackage[dvipsnames,svgnames,x11names,hyperref]{xcolor}
\definecolor{deepblue}{RGB}{0,35,102}
\definecolor{deepred}{RGB}{139,0,0}
\usepackage[table,xcdraw]{xcolor}

\usepackage{hyperref}
\hypersetup{
    colorlinks=true,
    linkcolor=RoyalPurple,
    urlcolor=RoyalPurple,
    citecolor=RoyalPurple
    }

\setlength{\belowdisplayskip}{0pt} \setlength{\belowdisplayshortskip}{0pt}
\setlength{\abovedisplayskip}{0pt} \setlength{\abovedisplayshortskip}{0pt}
\usepackage[font=footnotesize,labelfont=bf]{caption}

\usepackage{floatrow}
\usepackage{hyperref}
\newtheorem{proposition}{Proposition}
\newtheorem{lemma}{Lemma}
\newtheorem{definition}{Definition}
\newtheorem{remark}{Remark}

\usepackage[superscript,biblabel]{cite}
\linespread{1.0}


\title{Bispectral Signatures of Data: The Different Flavors of Bispectra and G-Bispectra}


\author{\vspace{-0.8cm}}
\date{}

\begin{document}


\maketitle
\tableofcontents

\paragraph{Triple correlation on R} The triple correlation of a function $f$ on $\mathbb{R}$ is:
$$
T(f)\left(t_1, t_2\right):=\int_{\mathbb{R}} f(x)^* f\left(x+t_1\right) f\left(x+t_2\right) d x .
$$
The triple correlation is translation invariant: if $g(x)=$ $f\left(x-x_0\right)$, then $T_g=T(f)$, as a consequence of the invariance of the integral.

We note that R is not included in any of the cases here: not finite, not compact, not homogeneous, etc - except the homogeneous x invariant space.

\section{Preliminaries}

Throughout the computations that follow, we use the following.

Properties of representations:
\begin{align*}
    \rho(gg') &= \rho(g) \rho(g')\\
    \rho(g^{-1}) = \rho(g)^\dagger
\end{align*}

Property of tensor product and matrix multiplication:
\begin{equation}
    (A \otimes B) . (C \otimes D) = (AC) \otimes (BD)
\end{equation}

\section{Bispectrum-only Summary}


\newpage
 \newgeometry{left=0cm,bottom=0cm}
\begin{table}[]
\vspace*{-20mm}
\centering
\rotatebox{90}{
\scriptsize{
\begin{tabular}{|l|l|l|}
\hline
\textbf{} &
  \textbf{Fourier} &
  Bispectrum \\ \hline
\textbf{$Z_n$} &
  \begin{tabular}[c]{@{}l@{}}$\mathcal{F}(f)_k$\\ $= \sum_{x=0}^{n-1} f(x).e^{-i2\pi k x / n}$\end{tabular} &
  \begin{tabular}[c]{@{}l@{}}$\beta(f)_{k_1, k_2} $\\ $= \hat{f}_{k_1}\hat{f}_{k_2}\hat{f}^{\dagger}_{k_1+k_2}$\end{tabular} \\ \hline
\textbf{$G_n$} &
  \begin{tabular}[c]{@{}l@{}}$\mathcal{F}(f)_\rho$\\ $=\frac{1}{|G|} \sum_{g \in G} f(g) \rho(g)^{\dagger}$\end{tabular} &
  \begin{tabular}[c]{@{}l@{}}$\beta(f)_{\rho_1, \rho_2}$\\ $=[\mathcal{F}(f)_{\rho_1} \otimes \mathcal{F}(f)_{\rho_2} ] C_{{\rho_1} {\rho_2}}\left[\bigoplus_{\rho \in \rho_1 \otimes \rho_2}\mathcal{F}(f)_\rho\right] C_{{\rho_1} {\rho_2}}^{\dagger}$\end{tabular} \\ \hline
\textbf{$G^c$} &
  same as $G$ &
  \begin{tabular}[c]{@{}l@{}}$\beta (f)_{\rho_1, \rho_2} $\\ $=  \mathcal{F}(f)_{\rho_1} \mathcal{F}(f)_{\rho_2}\mathcal{F}(f)_{\rho_1 \rho_2}^{\dagger}$\end{tabular} \\ \hline
\textbf{$S^2$} &
  \begin{tabular}[c]{@{}l@{}}$\mathcal{F}^v(f)_l $\\ $= \kappa_l \int_{R \in SO(3)} f^\dagger(T_R(x_0)). [\rho_l(R)]_0 dR$\end{tabular} &
  \begin{tabular}[c]{@{}l@{}}$\beta^v(f)_{l_1, l_2}[l]$\\ $:=\left[\mathcal{F}(f)_{l_1} \otimes \mathcal{F}(f)_{l_2} \right] C_{l_1l_2} \mathcal{F}_0(f)_l^{\dagger}$\end{tabular} \\ \hline
\textbf{} &
  \begin{tabular}[c]{@{}l@{}}$\mathcal{F}(f)_l $\\ $= \int_{R \in SO(3)} f^\dagger(T_R(x_0)). \rho_l(R) dR$\end{tabular} &
  \begin{tabular}[c]{@{}l@{}}$\beta(f)_{l_1l_2} $\\ $= [\mathcal{F}(f)_{l_1} \otimes \mathcal{F}(f)_{l_2}] C_{l_1,l_2} \Big[ \bigoplus_{l \in l_1 \otimes l_2} \mathcal{F}(f)_{l}^{\dagger} \Big] C^{\dagger}_{l_1,l_2}$\end{tabular} \\ \hline
\textbf{$H$} &
  \begin{tabular}[c]{@{}l@{}}$\mathcal{F}(f)_\rho $\\ $=\int_{g \in G} f\left(T_g(x_0)\right) \rho(g)dg$\end{tabular} &
  same as $G$ and $G_n$ \\ \hline
\textbf{$H^c$} &
  same as $H$ &
  same as $G^c$ \\ \hline
\end{tabular}
}
}
\caption{Summary}
\label{tab:eqns}
\end{table}

\section{Summary}

\newpage
 \newgeometry{left=0cm,bottom=0cm}
\begin{table}[]
\vspace*{-20mm}
\centering
\rotatebox{90}{
\scriptsize{
\begin{tabular}{|l|l|l|l|}
\hline
\textbf{} &
  \textbf{Fourier} &
  \textbf{Triple Correlation} &
  Bispectrum \\ \hline
\textbf{$Z_n$} &
  \begin{tabular}[c]{@{}l@{}}$\mathcal{F}(f)_k$\\ $= \sum_{x=0}^{n-1} f(x).e^{-i2\pi k x / n}$\end{tabular} &
  \begin{tabular}[c]{@{}l@{}}$T(f)_{x_1, x_2} $\\ $= \frac{1}{n} \sum_{x \in Z_n}f^*(x)f(x+x_1)f(x+x_2)$\end{tabular} &
  \begin{tabular}[c]{@{}l@{}}$\beta(f)_{k_1, k_2} $\\ $= \hat{f}_{k_1}\hat{f}_{k_2}\hat{f}^{\dagger}_{k_1+k_2}$\end{tabular} \\ \hline
\textbf{$G_n$} &
  \begin{tabular}[c]{@{}l@{}}$\mathcal{F}(f)_\rho $\\ $=\frac{1}{|G|} \sum_{g \in G} f(g) \rho(g)^{\dagger}$\end{tabular} &
  \begin{tabular}[c]{@{}l@{}}$T(f)_{g_1, g_2}$\\ $=\frac{1}{|G|} \sum_{g \in G} f^*(g) f\left(g g_1\right) f\left(g g_2\right)$\end{tabular} &
  \begin{tabular}[c]{@{}l@{}}$\beta(f)_{\rho_1, \rho_2}$\\ $=[\mathcal{F}(f)_{\rho_1} \otimes \mathcal{F}(f)_{\rho_2} ] C_{{\rho_1} {\rho_2}}\left[\bigoplus_{\rho \in \rho_1 \otimes \rho_2}\mathcal{F}(f)_\rho\right] C_{{\rho_1} {\rho_2}}^{\dagger}$\end{tabular} \\ \hline
\textbf{$G$} &
  \begin{tabular}[c]{@{}l@{}}$\mathcal{F}(f)_\rho $\\ $= \int_G f(g)\rho(g) dg$\end{tabular} &
  \begin{tabular}[c]{@{}l@{}}$T(f)_{g_1, g_2}$\\ $=\int_{g \in G} f^*(g) f\left(g g_1\right) f\left(g g_2\right)dg$\end{tabular} &
  same as $G_n$, i.e. as above \\ \hline
\textbf{$G^c$} &
  same as $G$ &
  same as $G$, i.e. as above &
  \begin{tabular}[c]{@{}l@{}}$\beta (f)_{\rho_1, \rho_2} $\\ $=  \mathcal{F}(f)_{\rho_1} \mathcal{F}(f)_{\rho_2}\mathcal{F}(f)_{\rho_1 \rho_2}^{\dagger}$\end{tabular} \\ \hline
\textbf{$S^2$} &
  \begin{tabular}[c]{@{}l@{}}$\mathcal{F}^v(f)_l $\\ $= \kappa_l \int_{R \in SO(3)} f^\dagger(T_R(x_0)). [\rho_l(R)]_0 dR$\end{tabular} &
  \begin{tabular}[c]{@{}l@{}}$T(f)_{R_1, R_2} $\\ $= \int_{R\in SO(3)} f^*(T_R(x_0)) f\left(T_{R R_1}(x_0)\right) f\left(T_{R R_2}(x_0)\right)dR$\end{tabular} &
  \begin{tabular}[c]{@{}l@{}}$\beta^v(f)_{l_1, l_2}[l]$\\ $:=\left[\mathcal{F}(f)_{l_1} \otimes \mathcal{F}(f)_{l_2} \right] C_{l_1l_2} \mathcal{F}_0(f)_l^{\dagger}$\end{tabular} \\ \hline
\textbf{} &
  \begin{tabular}[c]{@{}l@{}}$\mathcal{F}(f)_l $\\ $= \int_{R \in SO(3)} f^\dagger(T_R(x_0)). \rho_l(R) dR$\end{tabular} &
  Same as above, independent of Fourier transform. &
  \begin{tabular}[c]{@{}l@{}}$\beta(f)_{l_1l_2} $\\ $= [\mathcal{F}(f)_{l_1} \otimes \mathcal{F}(f)_{l_2}] C_{l_1,l_2} \Big[ \bigoplus_{l \in l_1 \otimes l_2} \mathcal{F}(f)_{l}^{\dagger} \Big] C^{\dagger}_{l_1,l_2}$\end{tabular} \\ \hline
\textbf{$H$} &
  \begin{tabular}[c]{@{}l@{}}$\mathcal{F}(f)_\rho $\\ $=\int_{g \in G} f\left(T_g(x_0)\right) \rho(g)dg$\end{tabular} &
  \begin{tabular}[c]{@{}l@{}}$T(f)_{g_1, g_2} $\\ $= \int_{g\in G} f^*(T_g(x_0)) f\left(T_{g g_1}(x_0)\right) f\left(T_{g g_2}(x_0)\right)dg$\end{tabular} &
  same as $G$ and $G_n$ \\ \hline
\textbf{$H^c$} &
  same as $H$ &
  same as $H$, i.e. as above &
  same as $G^c$ \\ \hline
\textbf{$D$} &
  \begin{tabular}[c]{@{}l@{}}$\mathcal{F}(f)_{nm} $\\ $= \frac{1}{a_{n m}} \int_0^{2 \pi} \int_0^1 f(r, \theta) J_m(2 \pi l_{n m} r) \exp(-i m \theta) r d r d\theta$\end{tabular} &
  \begin{tabular}[c]{@{}l@{}}$T(f)_{R_1, R_2} $\\ $:= \int_{R\in SO(2)} f^*(T_R(\theta_0), r_0) f\left(T_{R R_1}(\theta_0), r_0\right) f\left(T_{R R_2}(\theta_0), r_0\right)dR$\end{tabular} &
  \begin{tabular}[c]{@{}l@{}}$\beta(f)_{nmm'} $\\ $:= \mathcal{F}(f)_{nm}.\mathcal{F}(f)_{nm'}.\mathcal{F}(f)_{n(m+m')}^\dagger$\end{tabular} \\ \hline
\textbf{$B$} &
  \begin{tabular}[c]{@{}l@{}}$\mathcal{F}(f)_{l}$\\ $ \propto \int_0^{2 \pi}\int_0^{\pi} \int_0^1 f(r, \theta, \phi)^\dagger  \left( r^{\ell} Y_{\ell}(\theta, \phi)\right) r d r \cos(\theta) d\theta d\phi$\end{tabular} &
  \begin{tabular}[c]{@{}l@{}}$T(f)_{R_1, R_2} $\\ $:= \int_{R\in SO(3)} f^*(T_R(\theta_0, \phi_0), r_0) f\left(T_{R R_1}(\theta_0, \phi_0), r_0\right) f\left(T_{R R_2}(\theta_0, \phi_0), r_0\right)dR$\end{tabular} &
  $\beta^v(f)_{l_1, l_2}[l]$, same as $S^2$ \\ \hline
\textbf{$H \times R$} &
  \begin{tabular}[c]{@{}l@{}}$\mathcal{F}(f)_{\rho} $\\ $:= \int_{r \in R} \int_{g \in G} f\left(r, T_g(x_0)\right) \rho(g)dg. dr$\end{tabular} &
  \begin{tabular}[c]{@{}l@{}}$T(f)_{g_1, g_2} $\\ $:= \int_{g \in G} f^*(T_g(x_0), r_0) f\left(T_{gg_1}(x_0), r_0\right) f\left(T_{gg_2}(x_0), r_0\right)dg$\end{tabular} &
  Conjecture: $\beta^v$, same as $H$, \\ \hline
\end{tabular}
}
}
\caption{Summary}
\label{tab:eqns}
\end{table}
\restoregeometry
\newpage
\section{\color{deepblue}{Domain $\Omega = (\mathbb{Z}/n\mathbb{Z},~+)$ (cyclic translations)}}

\subsection{Fourier Transform: Definition}

The discrete Fourier transform of a signal $f$ on a space $\Omega$ is defined:

\vspace{-0.5cm}
\begin{equation}
  \hat{f}_k = \sum_{x=0}^{n-1}e^{-i2\pi k x / n} f(x).
\end{equation}
\vspace{-0.5cm}

From a group theoretic perspective, a Fourier transform is a projection of a signal onto the irreducible representations the group $(\mathbb{Z}/n\mathbb{Z}, +)$ of cyclic translations, indexed by integer frequencies $k$.

\subsection{Fourier Transform: Equivariance}
This gives it the property of \textit{equivariance} to translation in the input, commonly known as the time-shift and frequency-shift properties. Intuitively, this means that translating the input and then taking the Fourier transform is equivalent to taking the Fourier transform and then phase shifting the output. We expand the shift property of the Fourier transform here for illustration. Define the translate of $f$ by $z$ as $f^z(x) = f(x - z)$. Let $ h = x - z$. Then, The Fourier transform of $f^z(x)$ shifts as follows:

\vspace{-0.7cm}
\begin{eqnarray*}
  \hat{f}^z_k &=& \sum_{x=0}^{n-1}e^{-i2\pi k x / n} f(x - z) \\
%   \hat{f}^z(\omega) &=& \sum_{x=0}^{n-1} e^{\frac{-i2 \pi}{n}x\omega} f(x-z) \\
  \hat{f}^z_k &=& \sum_{h=0}^{n-1}e^{-i2\pi k (z+h) / n} f(h) \\
%   \hat{f}^z(\omega) &=& \sum_{h=0}^{n-1} e^{\frac{-i2 \pi}{n}(z+h)\omega} f(h) \\
  \hat{f}^z_k &=& \sum_{h=0}^{n-1} e^{-i2\pi k z / n} e^{-i2\pi k h / n} f(h) \\
%   \hat{f}^z(\omega) &=& \sum_{h=0}^{n-1} e^{\frac{-i2 \pi}{n}z\omega}e^{\frac{-i2 \pi}{n}h\omega} f(h) \\
  \hat{f}^z_k  &=&  e^{-i2\pi k z / n} \sum_{h=0}^{n-1} e^{-i2\pi k h / n} f(h) \\
%   \hat{f}^z(\omega) &=&  e^{\frac{-i2 \pi}{n}z\omega} \sum_{h=0}^{n-1} e^{\frac{-i2 \pi}{n}h\omega} f(h) \\
%   &=& e^{\frac{-i 2 \pi}{n}z \omega} \hat{f}(\omega)
% \end{eqnarray*}
  &=& e^{-i2\pi k z / n} \hat{f}_k.
\end{eqnarray*}
\vspace{-0.7cm}

That is, the shift is proportional to the translation $z$, by a factor of $e^{\frac{-i 2 \pi}{n} z k}$. Another way of thinking of the Fourier transform is as a homomorphism that maps group actions in the space domain (cyclic translation) to group actions in the frequency domain (complex multiplication).

\subsection{Steerable Bispectrum: After Convolution on Steerable Basis}

Question: can we compute the bispectrum more efficiently by formalizing the group convolution on a steerable basis? And recovering the output of the convolution directly in its Fourier form? We try to do it for the simplest bispectrum, i.e. the bispectrum on $\mathbb{Z}/n\mathbb{Z}$.

Consider a steerable basis for the discrete translations, i.e. the harmonics of the ring, where the ring is discretized into $\mathbb{Z} / n\mathbb{Z}$.

The basis $ x \mapsto Y_l(x) = e^{ilx}$ is steered by representations $\theta \mapsto \rho_l(\theta) = e^{-il\theta}$, since $Y_l(x - x_0) = \rho_l(-x_0)Y_l(x)$.

The filters of a convolutional layer are: $\Psi \in \mathbb{C}^{K \times n}$ for $K$ filters, given as their coefficients along the $n$ basis functions $Y_l$, $l=1, ..., n$. Let's consider only one filter for now, denoted $w$.
\begin{align*}
    w
    &= [\hat{w}^{1}, ..., \hat{w}^{n}] \in \mathbb{C}^n \\
    w(x)
    &= \sum_{l=1}^n \hat{w}^{(l)}Y_l(x), \quad \forall x \in \mathbb{Z} / n\mathbb{Z}
\end{align*}

The convolution between $w$ and a signal $f$ is, on the domain $\mathbb{Z} / n \mathbb{Z}$, is:
\begin{align*}
    \Theta(x_0) = (w \ast f) (x_0)
    &= \sum_{x \in \mathbb{Z}/n\mathbb{Z}} w(x_0-x).f(x) \\
    &= \sum_{x \in \mathbb{Z}/n\mathbb{Z}} \sum_{l=1}^n  \hat{w}^{(l)}Y_l(x_0-x).f(x) \\
    &= \sum_{x \in \mathbb{Z}/n\mathbb{Z}} \sum_{l=1}^n  \hat{w}^{(l)}\rho_l(-x). Y_l(x_0).f(x) \\
    &=\sum_{l=1}^n  \sum_{x \in \mathbb{Z}/n\mathbb{Z}} \hat{w}^{(l)}\rho_l(-x).f(x).Y_l(x_0) \\
    &=\sum_{l=1}^n \hat{\Theta}_l .Y_l(x_0),
    \quad \text{denoting:~} \hat{\Theta}_l \equiv \hat{w}^{(l)} \sum_{x \in \mathbb{Z}/n\mathbb{Z}} \rho_l(-x).f(x),
\end{align*}
with $\hat{\Theta}_l$ the coefficients of the output of the G-convolution, here the classic convolution, on the steerable basis.

We can directly compute the steerable bispectrum by combining the $\hat{\Theta}_l$ of the output.


\subsection{Triple Correlation: Definition}

$T(f)_{x_1, x_2} = \frac{1}{n} \sum_{x \in Z_n}f^*(x)f(x+x_1)f(x+x_2)$

\subsection{Triple Correlation: Symmetries}

\begin{proposition}[Symmetries]
The triple correlation of a signal has the following symmetry:
\begin{align*}
     (s1) \qquad T(f)_{x_1, x_2} &= T(f)_{x_2, x_1}
\end{align*}
The triple correlation of a real signal has the following additional symmetries:
\begin{align*}
     (s2) \qquad T(f)_{x_1, x_2} = T(f)_{-x_1, x_2-x_1} = T(f)_{x_2-x_1, -x_1} = T(f)_{x_1-x_2, -x_2} = T(f)_{-x_2, x_1-x_2}
\end{align*}
\end{proposition}

\begin{proof}
    We prove symmetry (s1), which relies on the fact that $f(x+x_2)$ and $f(x+x_1)$ are scalar values that commute:
    \begin{align*}
        T(f)_{x_2, x_1}
            = \frac{1}{n} \sum_{x \in Z_n} f^*(x)f(x+x_2)f(x+x_1)
            = \frac{1}{n} \sum_{x \in Z_n} f^*(x)f(x+x_1)f(x+x_2)
            = T(f)_{x_2, x_1}.
    \end{align*}
    We assume that the signal $f$ is real. We prove the first equality of symmetry (s2):
    \begin{align*}
        T(f)_{-x_1, x_2-x_1}
            &= \frac{1}{n} \sum_{x \in Z_n} f^*(x)f(x-x_1)f(x+x_2-x_1) \\
            &= \frac{1}{n} \sum_{x' \in Z_n} f^*(x'+x_1)f(x')f(x'+x_1+x_2-x_1)
                \quad \text{(with $x' = x - x_1$, i.e. $x = x'+x_1$)} \\
            &= \frac{1}{n} \sum_{x' \in Z_n} f^*(x'+x_1)f(x')f(x'+x_2)
                \quad \text{(because $Z_n$ commutative: $x_2 - x_1 = -x_1 + x_2$)} \\
            &= \frac{1}{n} \sum_{x \in Z_n} f(x+x_1)f(x)f(x+x_2)
                \quad \text{($f$ is a real signal, which equals its conjugate)}\\
            &= \frac{1}{n} \sum_{x \in Z_n} f(x)f(x+x_1)f(x+x_2)
                \quad \text{($f$ is takes on real values that commute)}\\
            &= \frac{1}{n} \sum_{x \in Z_n} f(x)^*f(x+x_1)f(x+x_2) \\
            &= T(f)_{x_2, x_1}
    \end{align*}
    The second equality of symmetry (s2) follows using (s1).
    The third and fourth equality of symmetry (s2) have the same proof.
\end{proof}

\begin{wrapfigure}{r}{0.25\textwidth}
  \centering
  \includegraphics[width=0.25\textwidth]{figs/tc_sectors.png}
  \caption{Sectors for the symmetries of the TC $\tau_1 = x_1$ and $\tau_2 = x_2$.}
  \label{fig:sectors}
\end{wrapfigure}
As a consequence, knowing the TC in any of the six sectors, I through VI, shown in Fig.~\ref{fig:sectors} would enable us to find the entire TC. These sectors include their boundaries so that, for example, sector $I$ is an infinite wedge bounded by the lines $\tau_2=0$ and $\tau_1=\tau_2, \tau_1 \geq 0$.


\subsection{Triple Correlation: Invariance}

\subsection{Triple Correlation: Continuity}

When $f_2$ is obtained by a group transformation of $f_1$, we have:
\begin{equation}
    f_2 = T_g[f_1] \Leftrightarrow T(f_2) = T(f_1).
\end{equation}

We explore what happens when $f_2$ is obtained from group-transformation of $f_1$ plus ``noise''. Do we have:
\begin{equation}
    f_2 = T_g[f_1] + \epsilon h \Rightarrow T(f_2) = T(f_1) + \epsilon T(h) \qquad?
\end{equation}


\begin{remark}
    We could also consider the case:
\begin{equation}
    f_2 = T_g[f_1] + T'_a[h] \Rightarrow T(f_2) = T(f_1) + ? \qquad?
\end{equation}
i.e. $f_2$ is obtained from group-transformation of $f_1$ plus a residual deformation by $a$ where $a$ is not the group defining the TC. For example, $a$ is from the general linear group.
\end{remark}

Consider the space $\mathcal{C}^0(G)$ of continuous signals $f$ defined on a compact group $G$. We equip this space with the uniform norm (also called sup norm): $|| f \|_{\infty}=\sup _{g \in G}|f(g)|$ for any $f \in \mathcal{C}^0(G)$. The operation ``computing the TC of $f$ at $g_1, g_2$'' can be expressed as:
$$
\begin{aligned}
& T: \mathcal{C}^0(G) \rightarrow \mathbb{R} \\
& f \rightarrow T(f)_{g_1, g2}.
\end{aligned}
$$

We prove the Lipschitz continuity of this operation.

\begin{proposition} There is the following bound on the TC of two signals $f_1, f_2 \in \mathcal{C}^0(G)$ for $G$ a locally compact group.
    \begin{align*}
 \|T\left(f_1\right)-T(f_2)\|_\infty
    \leq \text{vol}(G) \left( \|f_1\|_\infty^2 + \|f_1\|_\infty\|f_2\|_\infty + \|f_2\|_\infty^2\right) \left\|f_1-f_2\right\|_{\infty}.
\end{align*}
This shows the Lipschitz continuity of the TC:
    \begin{align*}
 \|T\left(f_1\right)-T(f_2)\|_\infty
    \leq K \left\|f_1-f_2\right\|_{\infty}.
\end{align*}
where $K = 3 \text{vol(G)}$ if $f_1, f_2$ are the outputs of a $G$-CNN with sigmoid non-linearity; or if $f_1, f_2$ are normalized signals that have max equal to 1. Alternatively, $K$ only depends on the weights of the $G$-CNN if $f_1, f_2$ are out to a $G$-CNN to which normalized signals have been sent.
\end{proposition}

\begin{proof}

Take two real signals $f_1$ and $f_2$ defined on a compact group $G$. Take two group elements $g_1, g_2 \in G$:
\begin{align*}
 \left|T\left(f_1\right)_{g_1,g_2}-T(f_2)_{g_1,g_2}\right|
    &=\left|\int_G f_1(g)f_1(gg_1)f_1(gg_2) dg - \int_G f_2(g)f_2(gg_1)f_2(gg_2)  dg\right| \\
    & \leq \int_G\left|f_1(g)f_1(gg_1)f_1(gg_2) -  f_2(g)f_2(gg_1)f(_2gg_2)\right| dg \\
    & \leq \text{vol}(G) ||F_3(f_1) - F_3(f_2)\|_{\infty},
\end{align*}
with notations: $F_3(f)(g) =  f(g)f(gg_1)f(gg_2)$ and $||F_3(f_1) - F_3(f_2)\|_{\infty} \equiv \sup_{g \in G} |F_3(f_1)(g) - F_3(f_2)(g)| $.

We will find an upper bound for $||F_3(f_1) - F_3(f_2)\|_{\infty}$. We have, for any $g$:
\begin{align*}
    |F_3(f_1)(g) - F_3(f_2)(g)|
        % &= |f_1(g)f_1(gg_1)f_1(gg_2) - f_2(g)f_2(gg_1)f_2(gg_2)|\\
        % &= |f_1(g)f_1(gg_1)f_1(gg_2) - f_1(g)f_1(gg_1)f_2(gg_2) \\
        % &~ + f_1(g)f_1(gg_1)f_2(gg_2) -  f_1(g)f_2(gg_1)f_2(gg_2) \\
        % &~ + f_1(g)f_2(gg_1)f_2(gg_2) -  f_2(g)f_2(gg_1)f_2(gg_2) |\\
        % &\leq |f_1(g)f_1(gg_1)f_1(gg_2) - f_1(g)f_1(gg_1)f_2(gg_2)| \\
        % & + |f_1(g)f_1(gg_1)f_2(gg_2) -  f_1(g)f_2(gg_1)f_2(gg_2)| \\
        % & + |f_1(g)f_2(gg_1)f_2(gg_2) -  f_2(g)f_2(gg_1)f_2(gg_2)| \\
        % & = |f_1(g)f_1(gg_1)(f_1(gg_2) -f_2(gg_2))|\\
        % & + |f_1(g)f_2(gg_2) (f_1(gg_1)-  f_2(gg_1))| \\
        % & + |f_2(gg_1)f_2(gg_2) (f_1(g)-  f_2(g))|\\
        % &\leq \left( |f_1(g)f_1(gg_1)| + |f_1(g)f_2(gg_2)| + |f_2(gg_1)f_2(gg_2)|\right) \left\|f_1-f_2\right\|_{\infty}\\
        &\leq \left( \|f_1\|_\infty^2 + \|f_1\|_\infty\|f_2\|_\infty + \|f_2\|_\infty^2\right) \left\|f_1-f_2\right\|_{\infty}
\end{align*}
using a telescopic sum, the triangle inequality and the definition of $\left\|.\right\|_{\infty}$.

The above inequality gives an upper bound that is valid for any $g \in G$. As the sup is defined as the lowest upper bound, we get:
\begin{equation}
    ||F_3(f_1) - F_3(f)\|_{\infty} \leq \left( \|f_1\|_\infty^2 + \|f_1\|_\infty\|f_2\|_\infty + \|f_2\|_\infty^2\right) \left\|f_1-f_2\right\|_{\infty}.
\end{equation}
and therefore:
\begin{align*}
 \left|T\left(f_1\right)_{g_1,g_2}-T(f_2)_{g_1,g_2}\right|
    \leq \text{vol}(G) \left( \|f_1\|_\infty^2 + \|f_1\|_\infty\|f_2\|_\infty + \|f_2\|_\infty^2\right) \left\|f_1-f_2\right\|_{\infty}.
\end{align*}
The above inequality gives an upper bound that is valid for any $g_1, g_2 \in G$. As the sup over $G \times G$ is defined as the lower upper bound, we get:
\begin{align*}
 \|T\left(f_1\right)-T(f_2)\|_\infty
    \leq \text{vol}(G) \left( \|f_1\|_\infty^2 + \|f_1\|_\infty\|f_2\|_\infty + \|f_2\|_\infty^2\right) \left\|f_1-f_2\right\|_{\infty}.
\end{align*}
This proves the Lipschitz continuity of the TC assuming that $f_1, f_2$ belong to a space of bounded signals.

For example, if the $G$-CNN uses a sigmoid as its nonlinearity, then for all $g \in G$: $|f_1(g)| \leq 1$ and $|f_1(g)| \leq 1$, and the above inequality proves the Lipschitz continuity of the TC with constant $K=3\text{vol}(G)$.

If the $G$-CNN has no nonlinearity, we recall its inputs are images with normalized intensity between 0 and 1. Thus, we can compute an upper bound of $\|f_1\|_\infty$ and $\|f_2\|_\infty$ that only depends on the (learned) weights of the $G$-CNN, and proves the Lipschitz continuity using this constant.
\end{proof}



\subsection{Bispectrum: Definition}

The Bispectrum is the Fourier transform of the triple correlation. Inserting the formula for the Fourier transform above, and simplifying (see details in compact groups proof):
\begin{align*}
    \beta(f)_{k_1, k_2}
        &=
    \frac{1}{n^3}\sum_{x_1 \in Z_n} \sum_{x_2 \in Z_n}
    T(f)(x_1, x_2)\left[
                \rho_1(x_1)^{\dagger} \otimes \rho_2(x_2)^{\dagger}
            \right]\\
    &=
    \frac{1}{n^3}\sum_{x_1 \in Z_n} \sum_{x_2 \in Z_n}
    \left(\sum_{x \in Z_n}
       f^*(x)
       f\left(x + x_1\right) f\left(x+ x_2\right)\right)
            \left[
                \rho_1(x_1)^{\dagger} \otimes \rho_2(x_2)^{\dagger}
            \right]\\
    &= \frac{1}{n}
        \sum_{x \in Z_n} f^*(x)
            [\mathcal{F}(f)_{\rho_1} \otimes \mathcal{F}(f)_{\rho_2}]
            \left[
                \rho_1(x) \otimes \rho_2(x)
            \right]\\
    &= \mathcal{F}(f)_{\rho_1} \otimes \mathcal{F}(f)_{\rho_2}
        \left[
            \frac{1}{n}
                \sum_{x \in Z_n} f(x)^* \rho_1(x) \otimes \rho_2(x)
        \right] \\
    &=
\mathcal{F}(f)_{\rho_1} \otimes \mathcal{F}(f)_{\rho_2} C_{{\rho_1} {\rho_2}}
    \left[
    \bigoplus_{\rho \in \rho_1 \otimes \rho_2}
        \mathcal{F}(f)_\rho
    \right] C_{{\rho_1} {\rho_2}}^{\dagger} .
\end{align*}
where each representation $\rho_j$ is associated with and integer $k_j$ such that: $\rho_1(x') = \exp(-i\frac{2\pi x'}{k_1})$ and $\rho_2(x') = \exp(-i\frac{2\pi x'}{k_2})$ for any $x' \in Z_n$.


For a pair of frequencies $k_1$ and $k_2$, the bispectrum $\beta$ is defined:

\begin{equation}
\beta_{k_1, k_2} = \hat{f}_{k_1}\hat{f}_{k_2}\hat{f}^{\dagger}_{k_1+k_2},
\end{equation}

with the summation $k_1+k_2$ taken modulo the dimension of the signal. Due to the conjugated third term, any absolute phase shifts resulting from translations in the input cancel out, leaving the bispectrum invariant.

\subsection{Bispectrum: Invariance}


Under translation, $\beta$ becomes:

\vspace{-1cm}
\begin{eqnarray*}
\beta^z_{k_1, k_2} &=& e^{-i 2 \pi k_1 z / n} \hat{f}_{k_1}  e^{-i 2 \pi k_2 z/ n} \hat{f}_{k_2} e^{i 2 \pi (k_1+k_2) z / n} \hat{f}^{\dagger}_{k_1 + k_2} \\
 &=& e^{-i 2 \pi k_1 z / n} e^{-i 2 \pi k_2 z/ n} e^{i 2 \pi k_1 z / n} e^{i 2 \pi k_2 z / n}\hat{f}_{k_1} \hat{f}_{k_2} \hat{f}^{\dagger}_{k_1 + k_2} \\
 &=& \hat{f}_{k_1} \hat{f}_{k_2} \hat{f}^{\dagger}_{k_1 + k_2}
\\ &=& \beta_{k_1, k_2}.
\end{eqnarray*}


\subsection{Bispectrum: Signal Reconstruction}

We answer the question: given the bispectrum $\beta(f)$, what is the minimum number of bispectral coefficients that we need to recover the signal $f$?

\begin{proposition}
    The $n+1$ coefficients $\beta(f)_{0, 0}, \beta(f)_{1, 0}, ...., \beta(f)_{1, n}$ are enough to recover the signal $f$.
\end{proposition}

\begin{lemma}
    We have the following relationships between bispectral coefficients:
    \begin{align*}
    \beta_{0, 0}
    &= \hat{f}_{0}\hat{f}_{0}\hat{f}^{\dagger}_{0}
    = \hat{f}_{0} | \hat{f}_{0} |^2 \quad \rightarrow |\beta_{0, 0}| = |\hat{f}_{0}|^3\\
    \beta_{1, 0}
        &= \hat{f}_{1} \hat{f}_{0} \hat{f}^{\dagger}_{1}
        =  \hat{f}_{0} | \hat{f}_{1} |^2, \\
    \beta_{k, 0}
        &= \hat{f}_{k} \hat{f}_{0} \hat{f}^{\dagger}_{k} =  \hat{f}_{0} |\hat{f}_{k}|^2 \\
    \beta_{1, k-1}
        &= \hat{f}_{1} \hat{f}_{k-1} \hat{f}^{\dagger}_{k} .
    \end{align*}
\end{lemma}

\begin{proof}
    We show that the $n+1$ coefficients $\beta(f)_{0, 0}, \beta(f)_{1, 0}, ...., \beta(f)_{1, n}$ are enough to recover the Fourier transform $\hat f$ of the signal $f$.

    The $\beta(f)_{0, 0}$ allows us to recover the DC component:
    \begin{equation}
            |\hat f (0)| = |\beta(f)_{0, 0}|^{1/3}, \quad \text{arg}(\hat f(0)) = \text{arg}(\beta_{0, 0}).
    \end{equation}

   We compute the initialization of the recurrence, i.e. the first frequency $\hat f (1)$:
   \begin{equation}
       | \hat f (1)| = \left(\frac{\beta_{1, 0}}{ \hat{f}_{0} }\right)^{1/2}, \quad  \text{arg}(\hat f(1)) = 0,
   \end{equation}
   We note that we could equally set $\text{arg}(\hat f(1)) = \phi$ for any $\phi \in [0, 2\pi)$. This indeterminacy is inevitable, since it is a reflection of translation invariance~\cite{kondor2008thesis}.

    Next, we compute the Fourier coefficients of $f$ iteratively, from the lowest frequency to the highest:
       \begin{equation}
       \hat f (k) = \sqrt{\frac{\beta_{1, k-1}}{ \hat{f}_{1} \hat{f}_{k-1} }}^\dagger.
   \end{equation}
\end{proof}


\section{\color{deepblue}{Domain $\Omega = (G,~+)$ (finite groups)}}

\subsection{Fourier Transform: Definition}

The (matrix-valued) Fourier transform on $G$ is defined by coefficients
$$
F(\omega)=\frac{1}{|G|} \sum_{g \in G} f(g) D_\omega(g)^{\dagger} .
$$


\subsection{Fourier Transform: Equivariance}

\subsection{Steerable Bispectrum: After G-Conv on Steerable Basis}

Consider the general case where $\{Y_l\}_l$ is a basis of steerable functions on a domain $\mathbb{\Omega}$ for a group $G$. That is the vector of stacked functions $Y$ is such that: $Y(g.x) = \rho(g)Y(x)$ for $\rho(g) \in \mathbb{C}^{L \times L}$ a representation of the group element $g$. We assume that such a basis exists.

We would like to compute the bispectrum to get a complete invariance for the action of the group $G$. We write the group convolution:
\begin{align*}
    \Theta(g_0) = (w \ast f)(g_0)
    &= \sum_{x \in \Omega}
        w(g_0^{-1}.x).f(x) \\
    &= \sum_{x \in \Omega}
        \sum_{l=1}^L \hat{w}^{l}Y_l(g_0^{-1}.x).f(x) \\
    &= \sum_{x \in \Omega}
        \hat{w}^T.Y(g_0^{-1}.x).f(x) \\
    &= \sum_{x \in \Omega}
        \left(\hat{w}^T.\rho(g_0)^{-1}.Y(x)\right).f(x) \\
    &= \sum_{x \in \Omega}
        \left(\hat{w}^T.\rho(g_0)^T.Y(x)\right).f(x),\quad \text{(assuming $\rho$ unitary)} \\
    &= \sum_{x \in \Omega}
        \left((\rho(g_0)\hat{w})^T.Y(x)\right).f(x),\\
    &= \sum_{x \in \Omega} \sum_{l=1}^L
        \left((\rho(g_0)\hat{w})_l.Y_l(x)\right).f(x),\\
    &=  \sum_{l=1}^L \left(
    \sum_{x \in \Omega}
        Y_l(x).f(x) \right)(\rho(g_0)\hat{w})_l,\\
    % &= (\rho(g_0)\hat{w})^T \sum_{x \in \Omega}
    %    Y(x).f(x),\quad\text{(with $Y(x) \in \mathbb{C}^L$ and $f(x) \in \mathbb{R}$)}.\\
\end{align*}
which ressembles a fourier decomposition:
\begin{equation}
    \Theta(x) = \sum_{l=1}^L \hat \Theta_l . Y_l(x)
\end{equation}

We would like to get the coefficients of $\hat{\Theta}$ on a steerable basis of the domain. However, the domain has changed. And we don't have a steerable basis of this new domain.

For the bispectrum, we would want the coefficients of $\Theta$ on a "basis" that is indexed by the irreducible representations of the group. Because the spectral coefficients are not indexed by $l$ anymore.

However, if we consider a case where the action is such that the output of the group convolution has the same domain as the input. Then we can use the same basis functions.


\subsection{Triple Correlation: Definition}

The triple correlation of a function $f$ defined on a finite group $G$ is:
$$
T(f)_{g_1, g_2}=\frac{1}{|G|} \sum_{g \in G} f^*(g) f\left(g g_1\right) f\left(g g_2\right) .
$$

\subsection{Triple Correlation: Invariance}
Note that $T(f)$ is a function defined on $G \times G$, and that, like the autocorrelation, the triple correlation is invariant under left translation: $T(f_1)=T(f_2)$ if $f_1(g)=f_2(h g)$.

\subsection{Triple Correlation: Symmetries}

\begin{proposition}[Symmetries]
The triple correlation of a signal has the following symmetry:
\begin{align*}
     (s1) \qquad T(f)_{g_1, g_2} &= T(f)_{g_2, g_1}
\end{align*}
If $G$ is commutative, the triple correlation of a real signal has the following additional symmetries:
\begin{align*}
     (s2) \qquad T(f)_{g_1, g_2}
     = T(f)_{g^{-1}_1, g_2g^{-1}_1}
     = T(f)_{g_2g^{-1}_1, g^{-1}_1}
     = T(f)_{g_1g^{-1}_2, g^{-1}_2}
     = T(f)_{g^{-1}_2, g_1g^{-1}_2}
\end{align*}
\end{proposition}


\begin{proof}
    We prove symmetry (s1), which relies on the fact that $f(gg_2)$ and $f(gg_1)$ are scalar values that commute:
    \begin{align*}
        T(f)_{g_2, g_1}
            = \frac{1}{|G|} \sum_{g \in G} f^*(g)f(gg_2)f(gg_1)
            = \frac{1}{|G|} \sum_{g \in G} f^*(g)f(gg_1)f(gg_2)
            = T(f)_{g_2, g_1}.
    \end{align*}
    We assume that the signal is real and that $G$ is commutative.
    We prove the first equality of symmetry (s2):
    \begin{align*}
        T(f)_{g^{-1}_1, g_2g^{-1}_1}
            &= \frac{1}{|G|} \sum_{g \in G} f^*(g)f(gg^{-1}_1)f(gg_2g^{-1}_1) \\
            &= \int_{g' \in G} f^*(g'g_1)f(g')f(g'g_1g_2g^{-1}_1)
                \quad \text{(with $g' = g - g_1$, i.e. $g = g'g_1$)} \\
            &= \int_{g' \in G} f^*(g'g_1)f(g')f(g'g_2)
                \quad \text{(because $G$ commutative: $g_2 g^{-1}_1 = g^{-1}_1 g_2$)} \\
            &= \frac{1}{|G|} \sum_{g \in G} f(gg_1)f(g)f(gg_2)
                \quad \text{($f$ is a real signal, which equals its conjugate)}\\
            &= \frac{1}{|G|} \sum_{g \in G} f(g)f(gg_1)f(gg_2)
                \quad \text{($f$ is takes on real values that commute)}\\
            &= \frac{1}{|G|} \sum_{g \in G} f(g)^*f(gg_1)f(gg_2) \\
            &= T(f)_{g_2, g_1}
    \end{align*}
    The second equality of symmetry (s2) follows using (s1).
    The third and fourth equality of symmetry (s2) have the same proof.
\end{proof}

\subsection{Bispectrum: Definition}

The bispectrum formula for arbitrary finite groups:
\begin{equation}
\beta(f)_{\rho_1, \rho_2}
=
\mathcal{F}(f)_{\rho_1} \otimes \mathcal{F}(f)_{\rho_2} C_{{\rho_1} {\rho_2}}
    \left[
    \bigoplus_{\rho \in \rho_1 \otimes \rho_2}
        \mathcal{F}(f)_\rho
    \right] C_{{\rho_1} {\rho_2}}^{\dagger} .
\end{equation}


\begin{proof}
We prove this by applying the Fourier transform to the triple correlation formula, following Karala. Note that the proof over compact (non-necessarily finite) groups will be similar (actually, this proof here can be deduced from the compact one).

The Fourier transformation of $T(f)$ requires Kronecker (tensor) products of the representation matrices $\left\{D_\omega\right\}_{\omega \in \Omega}$. The bispectrum is obtained by the formula:
\begin{equation}
    \beta(f)_{\rho_1, \rho_2}
    =
    \frac{1}{|G|^2} \sum_{g_1 \in G} \sum_{g_2 \in G}
        T(f)_{g_1, g_2} \rho_1(g_1)^{\dagger} \otimes \rho_2(g_2)^{\dagger}
\end{equation}
for all ${\rho_1}, {\rho_2}$. Note: Karala uses $|G|$ but it's probably $|G|^2$?

Inserting the formula for the Fourier transform above, and simplifying (see details in compact groups proof):
\begin{align*}
    \beta(f)_{\rho_1, \rho_2}
    &=
    \frac{1}{|G|^3}\sum_{g_1 \in G} \sum_{g_2 \in G} \sum_{g \in G}
       f^*(g)
       f\left(g g_1\right) f\left(g g_2\right)
            \left[
                \rho_1(g_1)^{\dagger} \otimes \rho_2(g_2)^{\dagger}
            \right] \\
    &= \frac{1}{|G|}
        \sum_{g \in G} f^*(g)
            [\mathcal{F}(f)_{\rho_1} \otimes \mathcal{F}(f)_{\rho_2}]
            \left[
                \rho_1(g) \otimes \rho_2(g)
            \right]\\
    &= \mathcal{F}(f)_{\rho_1} \otimes \mathcal{F}(f)_{\rho_2}
        \left[
            \frac{1}{|G|}
                \sum_{g \in G} f(g)^* \rho_1(g) \otimes \rho_2(g)
        \right] \\
    &=
\mathcal{F}(f)_{\rho_1} \otimes \mathcal{F}(f)_{\rho_2} C_{{\rho_1} {\rho_2}}
    \left[
    \bigoplus_{\rho \in \rho_1 \otimes \rho_2}
        \mathcal{F}(f)_\rho
    \right] C_{{\rho_1} {\rho_2}}^{\dagger} .
\end{align*}
where the last line is the Clebsh Gordan decomposition on $\rho_1 \times \rho_2$ and subsequent application of the definition of the Fourier transform.

The formula simplifies if $G$ is commutative. (TODO: proof?).

\end{proof}

\subsection{Bispectrum: Invariance}

\subsection{Bispectrum: Signal Reconstruction}

For the definition of the bispectrum, we take the convention with the transpose conjugate is on the first term, as in Eq 7.5 p 86 of Kondor thesis.

% SO2: C4 = Z/4Z -> OK

% Need the reduction for these three
% O2: D4

% D16 is the real goal: 8 rotations and 1 relfections

% SO3: Octahedral group with 24 elements
% https://en.wikipedia.org/wiki/Octahedral_symmetry
% O3: no chirality, with 48 elements
% --> Look into the e2cnn library / Weiler's work


We generalize to the non-commutative case.  We show that we recover the signal $f$ from its bispectral coefficients. We follow the same strategy as in the discrete, commutative case: we show that we can recover the Fourier transform of the signals for every irreducible representation $\rho$ of the group. We recover the signal itself by inverse Fourier transform. We first need to know ``how many'' irreducible representations the group has. For finite groups, this number is finite. We use the following

    \begin{proposition}
        The number of irreducible representations for a finite group $G$ is equal to the number of conjugacy classes. If $G$ is abelian, the number of irreducible representations is equal to $|G|$ and each representation is 1-dimensional.
    \end{proposition}

\begin{proof}
    We recall that the conjugacy class of an element $g \in G$ is defined with the conjugation operation defined by inner automorphisms:
    \begin{equation}
        C(g) = \{ hgh^{-1} ~|~ \forall h \in G\}.
    \end{equation}

    To prove the proposition in its entirety would require delving deep into the theory, but the core idea relies on the orthogonality relations of the character table of the group. The rows of the character table correspond to irreducible characters (which, over the complex numbers, can be associated with irreducible representations), while the columns correspond to conjugacy classes. Whether the group is commutative or not, this result holds. However, for commutative (or abelian) groups, every element is its own conjugacy class.
\end{proof}

We note that for Lie groups, this might not hold. For example, the group SO(3) has one irreducible representation in each odd dimension: it has an infinite number of irreps. We will deal with this case later.

We first show relationships between the bispectral coefficients and the Fourier coefficients. We denote $\rho_0$ the trivial representation, which is the representation that sends every group element to the identity matrix, i.e. to the scalar 1 by irreducibility of the $\rho$'s.

\begin{lemma}\label{lem:coefs}
    For $\rho_0$ the trivial representation of the group $G$, and $\rho$ any other irreducible representation of dimension $D$, the bispectral coefficients write:
    \begin{align*}
    \beta_{\rho_0, \rho_0}
    &= |\mathcal{F}(f)_{\rho_0}|^2 \mathcal{F}(f)_{\rho_0}
    \in \mathbb{C}\\
    \beta_{\rho, \rho_0}
    &= \mathcal{F}(f)_{\rho_0}^\dagger
    \mathcal{F}(f)_{\rho}^\dagger\mathcal{F}(f)_{\rho}
    \in \mathbb{C}^{D \times D}.
    \end{align*}
\end{lemma}

\begin{proof}
    Take $\rho_0$ the trivial representation, we know that $\rho \otimes \rho_0  = \rho$ such that the Clebsh-Gordan decomposition:
    \begin{equation}
        \rho \otimes \rho_0 = C_{{\rho} {\rho_0}}\left[
    \bigoplus_{\rho' \in \rho \otimes \rho_0}
        \rho'
    \right] C_{{\rho} {\rho_0}}^{\dagger} = C_{{\rho} {\rho_0}}
        \rho C_{{\rho} {\rho_0}}^{\dagger}
    \end{equation}
    gives: $C_{{\rho} {\rho_0}} C_{{\rho} {\rho_0}}^{\dagger}  = 1$ for any irreps $\rho$. In addition, for $\rho = \rho_0$, the Clebsh-Gordan matrices are scalars. Actually, in this case, the CJ matrix is just the identity matrix (according to chatGPT).

    We compute the bispectral coefficient:
    \begin{align*}
    \beta_{\rho_0, \rho_0}
    &=
    (\mathcal{F}(f)_{\rho_0} \otimes \mathcal{F}(f)_{\rho_0})^\dagger C_{{\rho_0} {\rho_0}}
    \left[
    \bigoplus_{\rho \in \rho_0 \otimes \rho_0}
        \mathcal{F}(f)_\rho
    \right] C_{{\rho_0} {\rho_0}}^{\dagger} \\
    &=
    (\mathcal{F}(f)_{\rho_0} \otimes \mathcal{F}(f)_{\rho_0})^\dagger C_{{\rho_0} {\rho_0}}
        \mathcal{F}(f)_{\rho_0}
     C_{{\rho_0} {\rho_0}}^{\dagger}
    \qquad \text{($\rho_0 \otimes \rho_0 = \rho_0$ which is irreducible)}
    \\
    &=
    (\mathcal{F}(f)_{\rho_0}^2)^\dagger C_{{\rho_0} {\rho_0}}
        \mathcal{F}(f)_{\rho_0}
     C_{{\rho_0} {\rho_0}}^{\dagger}
    \qquad \text{($\mathcal{F}(f)_{\rho_0}$ is a scalar for which tensor product is multiplication)}
    \\
    &=(\mathcal{F}(f)_{\rho_0}^\dagger)^2 \mathcal{F}(f)_{\rho_0}
    C_{{\rho_0} {\rho_0}} C_{{\rho_0} {\rho_0}}^{\dagger}
    \qquad\text{(scalars commute with other elements)}\\
    &=|\mathcal{F}(f)_{\rho_0}|^2 \mathcal{F}(f)_{\rho_0}
    \qquad\text{($C_{{\rho_0} {\rho_0}} C_{{\rho_0} {\rho_0}}^{\dagger}  = 1$).}
    \end{align*}

    Take any irreducible representation $\rho$ of dimension $D$, we have:
    \begin{align*}
        \beta_{\rho, \rho_0}
        &= (\mathcal{F}(f)_{\rho} \otimes \mathcal{F}(f)_{\rho_0})^\dagger C_{{\rho} {\rho_0}}
    \left[
    \bigoplus_{\rho \in \rho \otimes \rho_0}
        \mathcal{F}(f)_\rho
    \right] C_{{\rho} {\rho_0}}^{\dagger} \\
        &= \mathcal{F}(f)_{\rho_0}^\dagger \mathcal{F}(f)_{\rho}^\dagger C_{{\rho} {\rho_0}}
        \mathcal{F}(f)_{\rho} C_{{\rho} {\rho_0}}^{\dagger} \\
        &=\mathcal{F}(f)_{\rho_0}^\dagger \mathcal{F}(f)_{\rho}^\dagger\mathcal{F}(f)_{\rho} \qquad\text{(CJ coefficients for $\rho, \rho_0$ are identity matrices)}.
    \end{align*}
    \end{proof}

\begin{proposition}
    Consider a finite group $G$ that is non abelian. Consider the $L$ bispectral coefficients of an unknown signal $f$, where $L$ is a number computed from the Kroneck product table of the group. We can recover the Fourier coefficients of $f$ from only $L$ bispectral coefficients.
\end{proposition}


\begin{proof}
    We consider first Fourier coefficient (DC component), i.e the Fourier transform of the signal at the trivial representation $\rho_0$:
    \begin{equation}
    \mathcal{F}(f)_{\rho_0} = \hat{f}_{\rho_0} = \int_G f(g)\rho_0(g) dg = \int_G f(g) dg \in \mathbb{C}
     \end{equation}
    Using Lemma~\ref{lem:coefs}, we can recover it from the ``first'' bispectral coefficient, as:
    \begin{equation}
        | \mathcal{F}(f)_{\rho_0}| = \left(|\beta_{\rho_0, \rho_0}|\right)^{1/3},
        \quad
        \text{arg}(\mathcal{F}(f)_{\rho_0}) = \text{arg}(\beta_{\rho_0, \rho_0}),
    \end{equation}
    which mirrors the commutative case.

    We consider another Fourier coefficient $\mathcal{F}(f)_{\rho_1}$ for a representation $\rho_1$ of dimension $D$. The coefficient is a matrix in $\mathbb{C}^{D \times D}$. Using Lemma~\ref{lem:coefs}, we can recover it from a bispectral coefficient:
    \begin{equation}
        \mathcal{F}(f)_{\rho}^\dagger\mathcal{F}(f)_{\rho}
        = \frac{\beta_{\rho, \rho_0}}{\mathcal{F}(f)_{\rho_0}^\dagger} \in \mathbb{C}^{D \times D},
    \end{equation}
    where the denominator is a scalar in the equation above. The matrix $\mathcal{F}(f)_{\rho}^\dagger\mathcal{F}(f)_{\rho}$ is hermitian, and admits a square-root. We recover the Fourier coefficient as this square-root:
      \begin{equation}
        \mathcal{F}(f)_{\rho}'
        = \left(\frac{\beta_{\rho, \rho_0}}{\mathcal{F}(f)_{\rho_0}^\dagger}
        \right)^{1/2}.
    \end{equation}
    The square-root only recovers $\mathcal{F}(f)_{\rho}$ up to a matrix factor: this is an unidentifiability similar to the commutative case. Specifically, consider the SVD decomposition of $\mathcal{F}(f)_{\rho}$:
    \begin{align*}
        \mathcal{F}(f)_{\rho}
        &= U.\Sigma.V^\dagger\\
        \Rightarrow
         \mathcal{F}(f)_{\rho}^\dagger\mathcal{F}(f)_{\rho}
        &= (U.\Sigma.V^\dagger)^\dagger. U.\Sigma.V^\dagger
        = V \Sigma^2 V^\dagger\\
        \Rightarrow
        \mathcal{F}(f)_{\rho}'
        &= V \Sigma V^\dagger
    \end{align*}
    Thus, we have: $\mathcal{F}(f)_{\rho} = UV^\dagger.\mathcal{F}(f)_{\rho}' $ where $U, V$ are unitary matrices that come from the (unknown) SVD decomposition of the (unknown) $\mathcal{F}(f)_{\rho}$. By recovering $\mathcal{F}(f)_{\rho}$ as $\mathcal{F}(f)_{\rho}'$, we fix $UV^\dagger=I$ in the same way that we have fixed $\phi = 0$ (rotation of angle 0, i.e. the identity) in the commutative case.

    We seek to find the Fourier coefficient for another representation $\rho_2$. We operate as we would in the commutative case: we write the bispectral coefficient for $\rho_1, \rho_1$. (Note: we can recover it, up to unitary transformation. We only have to find what its UV is.) We denote $\mathcal{R}$ the (finite) set of irreducible representations of the group $G$.
    \begin{align*}
    \beta_{\rho_1, \rho_1}
    &=
    (\mathcal{F}(f)_{\rho_1} \otimes \mathcal{F}(f)_{\rho_1})^\dagger C_{{\rho_1} {\rho_1}}
    \left[
    \bigoplus_{\rho \in \rho_1 \otimes \rho_1}
        \mathcal{F}(f)_\rho
    \right] C_{{\rho_1} {\rho_1}}^{\dagger}\\
    &=
        (\mathcal{F}(f)_{\rho_1} \otimes \mathcal{F}(f)_{\rho_1})^\dagger
        C_{{\rho_1} {\rho_1}}
    \left[
    \bigoplus_{\rho \in \mathcal{R}}
        \mathcal{F}(f)_\rho^{n_{\rho, \rho_1}}
    \right] C_{{\rho_1} {\rho_1}}^{\dagger},
    \end{align*}
    where $n_{\rho, \rho_1}$ is the multiplicity of irreps $\rho$ in the decomposition of $\rho_1, \rho_1$. If we know the group $G$ a priori, then we can compute the decomposition of $\rho_1$ into its irreps and we know what this is.

    We get the equation:
    \begin{equation}
        \bigoplus_{\rho \in \mathcal{R}}
        \mathcal{F}(f)_\rho^{n_{\rho, \rho_1}}
    = C_{{\rho_1} {\rho_1}}^{-1}(\mathcal{F}(f)_{\rho_1} \otimes \mathcal{F}(f)_{\rho_1})^{-\dagger} \beta_{\rho_1, \rho_1} C_{{\rho_1} {\rho_1}}^{-\dagger},
    \end{equation}
    where we know everything on the right hand side.
    Therefore, every Fourier coefficient $\mathcal{F}(f)_\rho$ that appears in the decomposition of $\rho_1 \otimes \rho_1$ into irreducible irreps $\rho$ can be computed, by reading off the elements of the block diagonal matrix.  We get the $\mathcal{F}(f)_\rho$ for which $n_{\rho, \rho_1} \neq 0$.
    In the commutative case, this is how we get $\rho_2$.
    Here, we assume that we get at least one other Fourier coefficient, for a $\rho_2$ that we fix.

    We can then compute:
    \begin{align*}
    \beta_{\rho_1, \rho_2}
    &=
    (\mathcal{F}(f)_{\rho_1} \otimes \mathcal{F}(f)_{\rho_2})^\dagger C_{{\rho_1} {\rho_2}}
    \left[
    \bigoplus_{\rho \in \rho_1 \otimes \rho_2}
        \mathcal{F}(f)_\rho^{n_{\rho, \rho_2}}
    \right] C_{{\rho_1} {\rho_2}}^{\dagger}
    \end{align*}
    to get the equation:
        \begin{equation}
        \bigoplus_{\rho \in \mathcal{R}}
        \mathcal{F}(f)_\rho^{n_{\rho, \rho_2}}
    = C_{{\rho_1} {\rho_2}}^{-1}(\mathcal{F}(f)_{\rho_1} \otimes \mathcal{F}(f)_{\rho_2})^{-\dagger} \beta_{\rho_1, \rho_2} C_{{\rho_1} {\rho_2}}^{-\dagger},
    \end{equation}
    where everything on the right-hand side is known. We note that, for arbitrary groups and their representations, Clebsch–Gordan matrices are not known in general. Yet we assume that we know them for the groups that we consider, i.e. we have computed them numerically (possible, according to Wikipedia).


    Therefore, every Fourier coefficient $\mathcal{F}(f)_\rho$ that appears in the decomposition of $\rho_1 \otimes \rho_1$ into irreducible irreps $\rho$ can be computed, by reading off the elements of the block diagonal matrix. We get the $\mathcal{F}(f)_\rho$ for which $n_{\rho, \rho_2} \neq 0$.

    We iterate this procedure. The remaining question is: do we get all of the irreps in this way? If so: how many iterations do we need? The answer is that it depends on its CJ matrices.

    If we know the group a priori, we can compute the Kronecker product table for $G$, where at row $i$ and column $j$ we list all of the irreps that appear in the decomposition of $\rho_i \otimes \rho_j$.

    We want to compute $m_i$ i.e. the multiplicity $\rho_i$ in the decomposition $\tilde \rho = \rho_j \otimes \rho_k$. We present a simple computational procedure, inspired by Cohen \& Welling 2017, that can be used to determine that multiplicity.

    This procedure relies on the character $\chi_{\tilde \rho}(g)=\operatorname{Tr}(\tilde \rho(g))$ of the representation to be decomposed. The most important fact about characters is that the characters of irreps $\rho_i, \rho_j$ are orthogonal:
    \begin{equation}
    \left\langle\chi_{\rho_i,}, \chi_{\rho_k}\right\rangle \equiv \frac{1}{|G|} \sum_{h \in G} \chi_{\rho_i}(h) \chi_{\rho_k}(h)=\delta_{i k} .
    \end{equation}
    Furthermore, since the trace of a direct sum equals the sum of the traces (i.e. $\chi_{\rho \oplus \rho^{\prime}}=\chi_\rho+\chi_{\rho^{\prime}}$ ), and every representation $\rho$ is a direct sum of irreps, it follows that we can obtain the multiplicity of
    irrep $\rho_i$ in $\tilde \rho$ by computing the inner product with the $i$-th character:
    $$
    \left\langle\chi_{\tilde \rho}, \chi_{\rho_i}\right\rangle
    =
    \left\langle\chi_{\oplus_j m_j \rho_j}, \chi_{\rho_i}\right\rangle=\left\langle\sum_j m_j \chi_{\rho_j}, \chi_{\rho_i}\right\rangle
    =
    \sum_j m_j\left\langle\chi_{\rho_j}, \chi_{\rho_i}\right\rangle
    =m_i
    $$
    So the dot product of characters is all we need to determine the multiplicities that we need.


    \begin{table}[]
    \centering
    \begin{tabular}{|l|l|l|l|l|}
    \hline
    $\otimes$ & $\rho_0$       & $\rho_1$       & $\rho_2$       & $\rho_3$       \\ \hline
    $\rho_0$  & $(1, 0, 0, 0)$ & $(0, 1, 0, 0)$ & $(0, 0, 1, 0)$ & $(0, 0, 0, 1)$ \\ \hline
    $\rho_1$  & $(0, 1, 0, 0)$ & $(m_k^{11})_k$ & $(m_k^{12})_k$ & $(m_k^{13})_k$ \\ \hline
    $\rho_2$  & $(0, 0, 1, 0)$ & $(m_k^{21})_k$ & $(m_k^{22})_k$ & $(m_k^{23})_k$ \\ \hline
    $\rho_3$  & $(0, 0, 0, 1)$ & $(m_k^{31})_k$ & $(m_k^{32})_k$ & $(m_k^{33})_k$ \\ \hline
    \end{tabular}
    \caption{}
    \label{tab:my-table}
    \end{table}

    Once we have the Kronecker product table, then we can see how many bispectral coefficients we need to implement. This will depend on the group. If there are irreps missing from the row of $\rho_1$, then the answer is that we will need more than the tensor products of the form $\rho_1 \otimes \rho_j$ to get all of the coefficients. We might need to consider others of the form $\rho_1$.

    We consider the following procedure.
    \begin{itemize}
        \item List = $[\rho1]$
        \item Compute $\rho_1 \otimes \rho_1$.
        \item Take all the irreps that are not in List that appears in that decomposition.
        \item For every such irrep $\rho_2$:
        \begin{itemize}
            \item Compute $\rho_1 \otimes \rho_2$
            \item Append $\rho_2$ to List.
            \item Take all the irreps that are not in List that appears in that decomposition.
            \item For every such irrep $\rho_3$:
            \begin{itemize}
                \item Compute $\rho_1 \otimes \rho_3$
                \item Append $\rho_3$ to List.
                \item etc
            \end{itemize}
        \end{itemize}
    \end{itemize}

    We wonder whether this procedure gives us all the irreps of the group, and if so, in how many steps L. The number of steps is the number of bispectral coefficients needed.

    The proposed procedure essentially uses a breadth-first search algorithm on the space of irreducible representations, starting with $\rho_1$ and using the tensor product with $\rho_1$ as the mechanism to explore the space. Whether this procedure succeeds in including all the irreps in the List might on the group and its Kronecker (tensor) product table.

    However, for many groups, this procedure does indeed succeed in listing all the irreps. This is particularly true for groups where the tensor product of any two irreps spans the whole space of irreps. This is the case, for example, for SU(2) (the group of 2x2 unitary matrices with determinant 1).

\subsubsection{Example: Diedral group $D_4$}

    The diedral group $D_4$ is the group of symmetries of a square. $D_4$ has five irreps, e.g. shown in Cohen \& Welling 2017 (Steerable CNNs).

    \paragraph{Irreps and their matrices $\rho(g)$ for $g \in G$}

    \begin{figure}
        \centering
        \includegraphics[width=\textwidth]{figs/d4.png}
    \end{figure}

    \paragraph{Kronecker table} To get the Kronecker table, we compute each tensor-representation by:
    \begin{itemize}
        \item computing its matrices $\rho(g)$ on each group element, by definition of the tensor product of matrices
        \item computing the characters $\Xi(g) = \text{Tr}(\rho(g))$ for each group element
        \item computing the inner-product of the character $\Xi$ with each character of each irreps to get the multiplicity.
    \end{itemize}


\begin{table}[]
\centering
    \begin{tabular}{|l|l|l|l|l|l|}
    \hline
    $\otimes$ & $A_1$             & $A_2$             & $B_1$             & $B_2$             & $E$               \\ \hline
    $A_1$     & $(1, 0, 0, 0, 0)$ & $(0, 1, 0, 0, 0)$ & $(0, 0, 1, 0, 0)$ & $(0, 0, 0, 1, 0)$ & $(0, 0, 0, 0, 1)$ \\ \hline
    $A_2$     & $(0, 1, 0, 0, 0)$ & $(1, 0, 0, 0, 0)$ & $(0, 0, 0, 1, 0)$ & $(0, 0, 1, 0, 0)$ & $(0, 0, 0, 0, 1)$ \\ \hline
    $B_1$     & $(0, 0, 1, 0, 0)$ & $(0, 0, 0, 1, 0)$ & $(1, 0, 0, 0, 0)$ & $(0, 1, 0, 0, 0)$ & $(0, 0, 0, 0, 1)$ \\ \hline
    $B_2$     & $(0, 0, 0, 1, 0)$ & $(0, 0, 1, 0, 0)$ & $(0, 1, 0, 0, 0)$ & $(1, 0, 0, 0, 0)$ & $(0, 0, 0, 0, 1)$ \\ \hline
    $E$       & $(0, 0, 0, 0, 1)$ & $(0, 0, 0, 0, 1)$ & $(0, 0, 0, 0, 1)$ & $(0, 0, 0, 0, 1)$ & $(1, 1, 1, 1, 0)$ \\ \hline
    \end{tabular}
    \caption{Kronecker table for the diedral group $D_4$. We choose $\rho_0 = A_1$, $\rho_1 = E$.}
    \label{tab:my-table}
\end{table}


\paragraph{Procedure} We choose $\rho_0 = A_1$, $\rho_1 = E$.
\begin{itemize}
    \item List = [$\rho_0$, $\rho_1$]
    \item Compute irreps in the decomposition of $\rho_1 \otimes \rho_1$, that are not already in List. This gives: $A_2, B_1, B_2$.
    \item Stopping criterion reached: all irreps in List.
\end{itemize}

\paragraph{Computations after the $G$-Convolution}

We store the representations in a variable.

We compute every Fourier coefficient of $\Theta$, via:
\begin{equation}
    \mathcal{F}(\Theta)_\rho = \sum_{g \in G} \Theta(g).\rho(g)
\end{equation}
by taking a linear combination of the $\rho(g)$ matrices for each $\rho$ irreps. Note that $\mathcal{F}(\Theta)_\rho \in \mathbb{C}$ for every $\rho$ in $\{A_1, A_2, B_1, B_2\}$ and $\mathcal{F}(\Theta)_E \in \mathbb{C}^{2 \times 2}$.

We compute $\beta_{\rho_0, \rho_0} \in \mathbb{C}$.
\begin{equation}
    \beta_{\rho_0, \rho_0} = |\mathcal{F}(\Theta)_{\rho_0}|^2 \mathcal{F}(\Theta)_{\rho_0} \in \mathbb{C}
\end{equation}

We compute $\beta_{\rho_1, \rho_0} \in \mathbb{C}^{2 \times 2}$.
\begin{equation}
    \beta_{\rho_1, \rho_0} = \mathcal{F}(\Theta)_{\rho_1}^\dagger\mathcal{F}(\Theta)_{\rho_1} \mathcal{F}(\Theta)_{\rho_0}^* \in \mathbb{C}^{2 \times 2}
\end{equation}

We compute $\beta_{\rho_1, \rho_1} \in \mathbb{C}^{4 \times 4}$. We use the notation: $\mathcal{F} = \mathcal{F}(\Theta)_{\rho_1}$.
\begin{align*}
    \beta_{\rho_1, \rho_1}
    &= (\mathcal{F}(\Theta)_{\rho_1} \otimes \mathcal{F}(\Theta)_{\rho_1})^\dagger C_{{\rho_1} {\rho_1}}
    \left[
    \bigoplus_{\rho \in \rho_1 \otimes \rho_1=\{A_1, A_2, B_1, B_2\}}
        \mathcal{F}(\Theta)_\rho
    \right] C_{{\rho_1} {\rho_1}}^{\dagger} \\
    &= \left[\begin{array}{ll}
\mathcal{F}_{1,1}\left[\begin{array}{ll}
\mathcal{F}_{1,1} & \mathcal{F}_{1,2} \\
\mathcal{F}_{2,1} & \mathcal{F}_{2,2}
\end{array}\right] & \mathcal{F}_{1,2}\left[\begin{array}{ll}
\mathcal{F}_{1,1} & \mathcal{F}_{1,2} \\
\mathcal{F}_{2,1} & \mathcal{F}_{2,2}
\end{array}\right] \\
\mathcal{F}_{2,1}\left[\begin{array}{ll}
\mathcal{F}_{1,1} & \mathcal{F}_{1,2} \\
\mathcal{F}_{2,1} & \mathcal{F}_{2,2}
\end{array}\right] & \mathcal{F}_{2,2}\left[\begin{array}{ll}
\mathcal{F}_{1,1} & \mathcal{F}_{1,2} \\
\mathcal{F}_{2,1} & \mathcal{F}_{2,2}
\end{array}\right]
\end{array}\right]^\dagger C_{{\rho_1} {\rho_1}}
\begin{bmatrix}
     \mathcal{F}(\Theta)_{A_1} & 0 & 0 & 0\\
     0 & \mathcal{F}(\Theta)_{A_2} & 0 & 0\\
     0 & 0 & \mathcal{F}(\Theta)_{B_1} & 0\\
     0 & 0 & 0 & \mathcal{F}(\Theta)_{B_2}\\
\end{bmatrix}
C_{{\rho_1} {\rho_1}}^{\dagger}
\end{align*}

we need the CJ coefficients for the $\rho_1 \otimes \rho_1 = E \otimes E$ decomposition.

    % The Kronecker product table is not the character table of the group, but they are related through the orthogonality relations of the character values.

    % More specifically, the tensor product of two representations can be decomposed into irreducible representations, and the multiplicity of a particular irrep in this decomposition can be determined using the character values from the character table.



\end{proof}



\section{\color{deepblue}{Domain $\Omega = (G,~+)$ (compact groups)}}

\subsection{Fourier Transform: Definition}

Fourier analysis can be generalized to other compact groups, such as 2D and 3D rotation, by preserving this defining property of equivariance to the group action \cite{rudin1962fourier}.

\textbf{The (continuous) Fourier transform} (for commutative or non-commutative groups) is:
\begin{equation}
   \mathcal{F}(f)_\rho = \hat{f}_\rho = \int_G f(g)\rho(g) dg,
\end{equation}

defined for any $\rho$ being an irreducible representations, where $G$ denotes a compact group, and the integral uses the Haar measure $dg$ on $G$. Note that the definition of the Fourier transform is the same for non-commutative or commutative groups, the only requirement is for these groups to be compact.

\subsection{Fourier Transform: Equivariance}

Define the transformation of $f$ by $\tilde g$ as $(\tilde{g} \ast f)(g) = f({\tilde g}^{-1}\circ g)$. The Fourier transform undergoes a transformation by the same $\tilde g$, as follows, which shows its equivariance.

\vspace{-0.7cm}
\begin{align*}
  \mathcal{F}(\tilde{g} \ast f)_\rho
  &=  \int_{g \in G} (\tilde{g} \ast f)(g)\rho(g) dg \\
  &=  \int_{g \in G} f({\tilde g}^{-1}\circ g)\rho(g) dg \\
  &=  \int_{g' \in G} f(g')\rho(\tilde{g}\circ g') dg' \\
  &=  \int_{g' \in G} f(g')\rho(\tilde{g}).\rho(g') dg' \\
  &=  \rho(\tilde{g}) . \int_{g' \in G} f(g')\rho(g') dg' \\
  &=  \rho(\tilde{g}) .\mathcal{F}(f)_\rho
\end{align*}
\vspace{-0.7cm}

\subsection{Triple Correlation: Definition}

The triple correlation of a function $f$ defined on a compact group $G$ is:
$$
T(f)_{g_1, g_2}=\int_{g \in G} f^*(g) f\left(g g_1\right) f\left(g g_2\right)dg .
$$

\subsection{Triple Correlation: Invariance}
Note that $T(f)$ is a function defined on $G \times G$, and that, like the autocorrelation, the triple correlation is invariant under left translation: $T(f_1)=T(f_2)$ if $f_1(g)=f_2(h g)$.

\begin{proof}
Consider two real signals $f_1, f_2$ defined on a domain $\Omega$ equipped with the action of a group $G$, such that $f_2 = T_h[f_1]$ in the sense defined in the previous section.

We show that this implies that $T(f_1)_{g_1, g_2} = T(f_2)_{g_1, g_2}$. Taking $g_1, g_2 \in G$, we have:
\begin{align*}
    T(f_2)_{g_1, g_2}
        &= \int_{g \in G}
            f_2^*(g) f_2\left(gg_1\right) f_2\left(gg_2\right) dg  \\
        &= \int_{g \in G}
            T_h[f_1]^*(g) T_h[f_1]\left(gg_1\right) T_h[f_1]\left(gg_2\right) dg  \\
        &= \int_{g \in G}
            f_1^*(hg) f_1\left(hgg_1\right) f_1\left(hgg_2\right) dg  \\
        &= \int_{g \in G}
            f_1^*(g) f_1\left(gg_1\right) f_1\left(gg_2\right) dg \\
        &=  T(f_1)_{g_1, g_2}.
\end{align*}
where we use the change of variable $hg \rightarrow g$. This proves the invariance of the TC with respect to group actions on the signals.
\end{proof}

\subsection{Triple Correlation: Symmetries}

\begin{proposition}[Symmetries]
The triple correlation of a signal has the following symmetry:
\begin{align*}
     (s1) \qquad T(f)_{g_1, g_2} &= T(f)_{g_2, g_1}
\end{align*}
If $G$ is commutative, and $f$ takes on real values, the triple correlation has the following additional symmetries:
\begin{align*}
     (s2) \qquad T(f)_{g_1, g_2}
     = T(f)_{g^{-1}_1, g_2g^{-1}_1}
     = T(f)_{g_2g^{-1}_1, g^{-1}_1}
     = T(f)_{g_1g^{-1}_2, g^{-1}_2}
     = T(f)_{g^{-1}_2, g_1g^{-1}_2}
\end{align*}
\end{proposition}


\begin{proof}
    We prove symmetry (s1), which relies on the fact that $f(gg_2)$ and $f(gg_1)$ are scalar values that commute:
    \begin{align*}
        T(f)_{g_2, g_1}
            = \int_{g \in G}  f^*(g)f(gg_2)f(gg_1) dg
            = \int_{g \in G}  f^*(g)f(gg_1)f(gg_2) dg
            = T(f)_{g_2, g_1}.
    \end{align*}
    We assume that the signal is real and that $G$ is commutative.
    We prove the first equality of symmetry (s2):
    \begin{align*}
        T(f)_{g^{-1}_1, g_2g^{-1}_1}
            &= \int_{g \in G} f^*(g)f(gg^{-1}_1)f(gg_2g^{-1}_1) dg\\
            &= \int_{g' \in G} f^*(g'g_1)f(g')f(g'g_1g_2g^{-1}_1) dg
                \quad \text{(with $g' = g - g_1$, i.e. $g = g'g_1$)} \\
            &= \int_{g' \in G} f^*(g'g_1)f(g')f(g'g_2) dg
                \quad \text{(because $G$ commutative: $g_2 g^{-1}_1 = g^{-1}_1 g_2$)} \\
            &= \int_{g \in G}  f(gg_1)f(g)f(gg_2) dg
                \quad \text{($f$ is a real signal, which equals its conjugate)}\\
            &= \int_{g \in G}  f(g)f(gg_1)f(gg_2) dg
                \quad \text{($f$ is takes on real values that commute)}\\
            &= \int_{g \in G}  f(g)^*f(gg_1)f(gg_2) dg\\
            &= T(f)_{g_2, g_1}
    \end{align*}
    The second equality of symmetry (s2) follows using (s1).
    The third and fourth equality of symmetry (s2) have the same proof.
\end{proof}

\subsection{Bispectrum: Definition}

The bispectrum can be defined more generally for signals defined over a space that is homogeneous for a compact group \cite{kakarala2012bispectrum}.

\textbf{The bispectrum for commutative groups} is defined over a pair of irreducible representations $\rho_1, \rho_2$:

\begin{equation}
    \beta (f)_{\rho_1, \rho_2} = \hat{f}_{\rho_1}\hat{f}_{\rho_2}\hat{f}_{\rho_1 \rho_2}^{\dagger} \qquad \in \mathbb{C}
    \label{eqn:group-bs}
\end{equation}
The irreducible representations (irreps) of commutative groups map to scalar values.

\textbf{The bispectrum for non-commutative groups} is defined over a pair of irreducible representations $\rho_1, \rho_2$:

\begin{align*}
   \beta (f)_{\rho_1, \rho_2}
   &= [\mathcal{F}(f)_{\rho_1} \otimes \mathcal{F}(f)_{\rho_2}] C_{\rho_1,\rho_2} \Big[ \bigoplus_{\rho \in \rho_1 \otimes \rho_2} \mathcal{F}(f)_{\rho}^{\dagger} \Big] C^{\dagger}_{\rho_1,\rho_2}\\
   &= [\hat{f}{\rho_1} \otimes \hat{f}{\rho_2}] C_{\rho_1,\rho_2} \Big[ \bigoplus_{\rho \in \rho_1 \otimes \rho_2} \hat{f}_{\rho}^{\dagger} \Big] C^{\dagger}_{\rho_1,\rho_2} \qquad \in \mathbb{C}^{D_1 D_2 \times D_1 D_2}
    \label{eqn:non-commutative-bs}
\end{align*}

The irreps of \textit{non-commutative} groups map to matrices. In this definition, $C$ is a Clebsch-Gordan matrix, $\otimes$ is the tensor product, and $\oplus$ is a direct sum over irreps. The unitary Clebsch-Gordan matrix $C$ is analytically defined for each pair of representations as:

\begin{equation}
    (\rho_1 \otimes \rho_2)(g) = C_{\rho_1,\rho_2}^\dagger \Big[ \bigoplus_{\rho \in \rho_1 \otimes \rho_2} \rho(g) \Big] C_{\rho_1,\rho_2}.
    \label{eqn:non-commutative-bs}
\end{equation}

\begin{proof}
We prove this by applying the Fourier transform to the triple correlation formula, generalizing the proof of Karakala to compact groups.

The Fourier transformation of $T(f)$ requires tensor products of the representation matrices, since it is defined on the product group.

We denote $F = \mathcal{F}(f)$ for simplicity of notations.

The bispectrum is obtained by the formula:
$$
\beta(f)_{\rho_1, \rho_2}
= \int_{g_1\in G} \int_{g_2\in G}
    T(f)_{g_1, g_2}
    \left[
        \rho_1\left(g_1\right)^{\dagger} \otimes \rho_2\left(g_2\right)^{\dagger}
    \right]dg_1dg_2
$$
for all ${\rho_1}, {\rho_2}$. Inserting the definition of the triple correlation, and simplifying with a change of variable $g_1 = gg_1$, $g_2 = gg_2$:
\begin{align*}
\beta(f)_{\rho_1, \rho_2}
=&
\int_{g_1\in G} \int_{g_2\in G}
    \int_{g \in G} f^*(g) f\left(g g_1\right) f\left(g g_2\right)
    \left[
        \rho_1\left(g_1\right)^{\dagger} \otimes \rho_2\left(g_2\right)^{\dagger}
    \right]
    dg dg_1dg_2 \\
=& \int_{g \in G}
    f^*(g)
    \int_{g_1\in G} \int_{g_2\in G}
    f\left(g g_1\right) f\left(g g_2\right)
    \left[
        \rho_1\left(g_1\right)^{\dagger} \otimes \rho_2\left(g_2\right)^{\dagger}
    \right]
    dg_1dg_2
    dg \\
=& \int_{g \in G}
    f^*(g)
    \int_{g_1\in G} \int_{g_2\in G}
    f\left(g_1\right) f\left(g_2\right)
    \left[
        \rho_1\left(g^{-1}g_1\right)^{\dagger} \otimes \rho_2\left(g^{-1}g_2\right)^{\dagger}
    \right]
    dg_1dg_2
    dg \\
=& \int_{g \in G}
    f^*(g)
    \int_{g_1\in G} \int_{g_2\in G}
    f\left(g_1\right) f\left(g_2\right)
    \left[
        \rho_1(g_1)^{\dagger}\rho_1(g^{-1})^{\dagger} \otimes \rho_2(g_2)^{\dagger}\rho_2(g^{-1})^\dagger
    \right]
    dg_1dg_2
    dg \\
=& \int_{g \in G}
    f^*(g)
    \int_{g_1\in G} \int_{g_2\in G}
    f\left(g_1\right) f\left(g_2\right)
    \left[
        \rho_1(g_1)^{\dagger}\rho_1(g) \otimes \rho_2(g_2)^{\dagger}\rho_2(g)
    \right]
    dg_1dg_2
    dg \\
=& \int_{g \in G}
    f^*(g)
    \int_{g_1\in G} \int_{g_2\in G}
    f\left(g_1\right) f\left(g_2\right)
        \left[\rho_1(g_1)^{\dagger}\otimes  \rho_2(g_2)^{\dagger}\right] . \left[\rho_1(g) \otimes \rho_2(g)\right]
    dg_1dg_2
    dg \\
&= \int_{g \in G}
    f^*(g)
    [\mathcal{F}(f)_{\rho_1} \otimes \mathcal{F}(f)_{\rho_2}]\left[\rho_1(g) \otimes \rho_2(g)\right]
    dg\\
&=
\mathcal{F}(f)_{\rho_1} \otimes \mathcal{F}(f)_{\rho_2}
    \left[
        \int_{g \in G} f(g)^* \rho_1(g) \otimes \rho_2(g) dg
    \right].\\
&=
\mathcal{F}(f)_{\rho_1} \otimes \mathcal{F}(f)_{\rho_2}
    \left[
        \int_{g \in G} f(g)^*  \left( C_{\rho_1 \rho_2} \bigoplus_{\rho \in \rho_1 \otimes \rho_2} \rho(g) C_{\rho_1 \rho_2}^{\dagger}\right) dg
    \right].\\
& =
\mathcal{F}(f)_{\rho_1} \otimes \mathcal{F}(f)_{\rho_2} C_{\rho_1 \rho_2} \left[ \bigoplus_{\rho \in \rho_1 \otimes \rho_2} \mathcal{F}(f)_{\rho} \right]  C_{\rho_1 \rho_2}^{\dagger}
\end{align*}

where the last two lines uses the Clebsh-Gordan formula to reduce the tensor product of representations.


The formula simplifies to (1) if $G$ is Abelian, in which case the fundamental theorem of finite Abelian groups shows that $G$ is a direct product of cyclic groups, each of which is represented by the complex exponentials $x \mapsto e^{j \omega x}$.

\end{proof}


\subsection{Bispectrum: Invariance}

Define the transformation of $f$ by $\tilde g$ as $(\tilde{g} \ast f)(x) = f({\tilde g}^{-1}\circ g)$. Under this group action, the bispectrum becomes:
\begin{align*}
   \beta(\tilde{g} \ast f)_{\rho_1, \rho_2}
   &= \Big[\mathcal{F}(\tilde{g} \ast f)_{\rho_1} \otimes \mathcal{F}(\tilde{g} \ast f)_{\rho_2})\Big] .C_{\rho_1,\rho_2} \Big[ \bigoplus_{\rho \in \rho_1 \otimes \rho_2} \mathcal{F}(\tilde{g} \ast f)_{\rho}^{\dagger} \Big]. C^{\dagger}_{\rho_1,\rho_2}
   \quad \text{(bispectrum definition)}\\
   &= \Big[(\rho_1(\tilde{g}).\mathcal{F}(f)_{\rho_1}) \otimes (\rho_2(\tilde{g}). \mathcal{F}(f)_{\rho_2})\Big] .C_{\rho_1,\rho_2} \Big[ \bigoplus_{\rho \in \rho_1 \otimes \rho_2} \rho(\tilde{g})^\dagger .\mathcal{F}(f)_{\rho}^{\dagger} \Big] .C^{\dagger}_{\rho_1,\rho_2}
   \quad \text{(equivariance of $\mathcal{F}$)}\\
  &=
  \Big[\rho_1(\tilde{g}) \otimes \rho_2(\tilde{g})\Big] .
  \Big[\mathcal{F}(f)_{\rho_1} \otimes \mathcal{F}(f)_{\rho_2})\Big] .
  C_{\rho_1,\rho_2}
  \Big[ \bigoplus_{\rho \in \rho_1 \otimes \rho_2} \rho(\tilde{g})^\dagger .\mathcal{F}(f)_{\rho}^{\dagger} \Big] .C^{\dagger}_{\rho_1,\rho_2}
  \quad \text{(*)}\\
    &=
  \Big[\rho_1(\tilde{g}) \otimes \rho_2(\tilde{g})\Big] .
  \Big[\mathcal{F}(f)_{\rho_1} \otimes \mathcal{F}(f)_{\rho_2})\Big] .
  C_{\rho_1,\rho_2} .
  \Big[ \bigoplus_{\rho \in \rho_1 \otimes \rho_2} \rho(\tilde{g})^\dagger) \Big].\Big[\bigoplus_{\rho \in \rho_1 \otimes \rho_2} \mathcal{F}(f)_{\rho}^{\dagger} \Big]. C^{\dagger}_{\rho_1,\rho_2}
  \quad \text{($\oplus$ definition)}\\
  &=
  \Big[\rho_1(\tilde{g}) \otimes \rho_2(\tilde{g})\Big] .
  \Big[\mathcal{F}(f)_{\rho_1} \otimes \mathcal{F}(f)_{\rho_2})\Big] .
  C_{\rho_1,\rho_2}. \Big[ \bigoplus_{\rho \in \rho_1 \otimes \rho_2} \rho(\tilde{g})^\dagger) \Big] .C^{\dagger}_{\rho_1,\rho_2}.C_{\rho_1,\rho_2}.\Big[\bigoplus_{\rho \in \rho_1 \otimes \rho_2} \mathcal{F}(f)_{\rho}^{\dagger} \Big] .C^{\dagger}_{\rho_1,\rho_2}
  \quad \text{(unitary)}\\
  &=
  \Big[\rho_1(\tilde{g}) \otimes \rho_2(\tilde{g})\Big]
  .
  \Big[\mathcal{F}(f)_{\rho_1} \otimes \mathcal{F}(f)_{\rho_2})\Big] .
  \Big[ C_{\rho_1,\rho_2}. \bigoplus_{\rho \in \rho_1 \otimes \rho_2} \rho(\tilde{g})^\dagger) .C^{\dagger}_{\rho_1,\rho_2}\Big].
  C_{\rho_1,\rho_2}.\Big[\bigoplus_{\rho \in \rho_1 \otimes \rho_2} \mathcal{F}(f)_{\rho}^{\dagger} \Big] .C^{\dagger}_{\rho_1,\rho_2}\\
  &=
  \Big[\rho_1(\tilde{g}) \otimes \rho_2(\tilde{g})\Big] .\Big[\mathcal{F}(f)_{\rho_1} \otimes \mathcal{F}(f)_{\rho_2})\Big] .\Big[\rho_1(\tilde{g}) \otimes \rho_2(\tilde{g})\Big]^\dagger. C_{\rho_1,\rho_2}.\Big[\bigoplus_{\rho \in \rho_1 \otimes \rho_2} \mathcal{F}(f)_{\rho}^{\dagger} \Big] .C^{\dagger}_{\rho_1,\rho_2}\\
  &=
  \Big[(\rho_1(\tilde{g}).\rho_1(\tilde{g})^\dagger.\mathcal{F}(f)_{\rho_1}) \otimes (\rho_2(\tilde{g}).\rho_2(\tilde{g})^\dagger. \mathcal{F}(f)_{\rho_2})\Big].  C_{\rho_1,\rho_2}.\Big[\bigoplus_{\rho \in \rho_1 \otimes \rho_2} \mathcal{F}(f)_{\rho}^{\dagger} \Big]. C^{\dagger}_{\rho_1,\rho_2}\\
  &=
  \Big[\mathcal{F}(f)_{\rho_1} \otimes\mathcal{F}(f)_{\rho_2}\Big] . C_{\rho_1,\rho_2}.\Big[\bigoplus_{\rho \in \rho_1 \otimes \rho_2} \mathcal{F}(f)_{\rho}^{\dagger} \Big]. C^{\dagger}_{\rho_1,\rho_2}\qquad \text{(unitary representations)}\\
  &=
  \beta(f)_{\rho_1, \rho_2} \qquad \text{(bispectrum definition)}
\end{align*}
where we use the linear algebra identity $(*)$: $(A.B) \otimes (C.D) = (A\otimes C) . (B \otimes D)$ for $\otimes$ the tensor product and $.$ the matrix multiplication.

\section{\color{purple}{Domain $\Omega = (S^2,~SO(3),~\ast)$ (spherical rotations)}}

We move towards homogeneous space for a group, and first consider an example before writing the general theory. This example is the spherical harmonic transform, which is a generalization of the Fourier transform to the non-commutative group of spherical rotations $SO(3)$ acting on the 2-sphere.

\subsection{Fourier Transform: Definition}
Functions defined over a sphere $S^2$ can be expanded on the spherical harmonics as:
\begin{equation}
    f(\theta, \phi)=\sum_{l=0}^{\infty} \sum_{m=-l}^l \widehat{f}_{l, m} Y_l^m(\theta, \phi)
\end{equation}
where $Y_l^m(\theta, \phi)$ as the spherical harmonics of coefficients $l, m$ given by $Y_l^m(\theta, \phi)=k_{lm} P_l^m(\cos \theta) e^{i m \phi}$
where $l=0,1,2, \ldots ; \quad m=[-l, l]$, $P_l^m$ are the associated Legendre polynomials, and $k_{lm}-\sqrt{\frac{2 l+1}{4 \pi} \frac{(l-m) !}{(l+m) !}}$ the normalization constant.

We recall that the spherical harmonics are the eigenfunctions of the Laplace operator on $S_2$ (with eigenvalue $-l^2$ ), and they form an orthonormal basis for $L_2\left(S_2\right)$.

\textbf{The spherical harmonics transform} of a spherical function $f$ is given by:
\begin{align*}
    \widehat{f}_{l, m}
    &= <f, Y_l^m>   \qquad \in \mathbb{C} \\
    &= \int_0^\pi \int_0^{2 \pi} f^\dagger(\theta, \phi) Y_l^m(\theta, \phi) \cos \theta d \phi d \theta \\
    &= \int_0^\pi \int_0^{2 \pi} f^\dagger(\theta, \phi) k_{lm} P_l^m(\cos \theta) e^{i m \phi} \cos \theta d \phi d \theta \\
    &=  k_{lm}   \int_0^\pi \Big[ \int_0^{2 \pi} f^\dagger(\theta, \phi)e^{i m \phi}d \phi\Big] P_l^m(\cos \theta)\cos \theta  d \theta
\end{align*}
where the last line shows that the spherical harmonic transform can be reduced to a regular Fourier transform in the longitudinal coordinate $\phi$ followed by a projection onto the associated Legendre functions. We note that there exists an expression for the discrete Fourier transform associated to that case.


We bring this expression into an equation amenable to a definition of generalized Fourier transform, i.e. only keeping the index $l$ that indexes the irreducible representations of $SO(3)$. Thus, we get vector-valued Fourier coefficients.

\textbf{The (vector) Fourier transform} of a spherical function $f$, where the subscript $0$ will become clear at the end of this computation, is:
\begin{align*}
  \mathcal{F}_0(f)_l = \widehat{f}_{l}
    &= <f, Y_l> \\
    &= \int_0^\pi \int_0^{2 \pi} f^\dagger(\theta, \phi) Y_l(\theta, \phi) \cos \theta d \phi d \theta   \qquad \in \mathbb{C}^{2l+1} \\
    &= \sqrt{\frac{2l+1}{4\pi}}\int_0^\pi \int_0^{2 \pi} f^\dagger(\theta, \phi). [\rho_l(\phi, \theta, \gamma)]_0 \cos \theta d \phi d \theta
    \qquad \text{(relation between $Y$'s and Wigner matrices, Karakala)}\\
    &= \frac{1}{\pi}\sqrt{\frac{2l+1}{4\pi}}\int_0^\pi\int_0^\pi \int_0^{2 \pi} f^\dagger(\theta, \phi). [\rho_l(\phi, \theta, \gamma)]_0 \cos \theta d \phi d \theta d\gamma \\
    &= \kappa_l \int_{R \in SO(3)} f^\dagger(T_R(x_0)). [\rho_l(R)]_0 dR \qquad \text{(Haar measure in Euler angles)}
\end{align*}
where $[\rho_l(R)^*]_0$ is the first column of the Wigner matrix, and $R$ is a rotation bringing the north pole $x_0$ into $(\theta, \phi)$, i.e. the rotation of Euler angles $(\alpha, \beta, \gamma)$ where $\alpha=\phi$ and $\beta=\theta$, i.e. convention of Euler angles, where $\alpha$  is a longitudinal angle and
$\beta$  is a colatitudinal angle (spherical polar angles in the physical definition of such angles). We use $\kappa_l$ to denote a normalization constant.

We note that this is \textbf{not equal} to the Fourier transform defined on an homogeneous space, as only one column is extracted: we have a vector-valued Fourier transform, instead of a matrix-valued Fourier transform, which is more computationally effective. This raises the question on whether every Fourier transform on homogeneous spaces can be reduced into a vector-valued transform in a more general way.

\textbf{The Fourier transform} of a spherical function $f$ is indeed:
\begin{align*}
  \mathcal{F}(f)_l
    = \int_{R \in SO(3)} f^\dagger(T_R(x_0)). \rho_l(R) dR.
\end{align*}

\subsection{Fourier Transform: Equivariance}


Define the transformation of $f$ by $\tilde R$ as $(\tilde{R} \ast f)(\theta, \phi) = f(\tilde R (\theta, \phi)) = f(T_{\tilde R}^{-1}.u)$. The (vector) Fourier transform undergoes a transformation by the same $\tilde R$, as follows, which shows its equivariance.

\begin{align*}
    \mathcal{F}(\tilde{R} \ast f)_{l}
    &=\kappa_l \int_{R \in SO(3)} (\tilde{R} \ast f)^\dagger(T_R(x_0)). [\rho_l(R)]_0 dR\\
    &=\kappa_l \int_{R \in SO(3)}
        f^\dagger(T_{\tilde R}^{-1} \circ T_R(x_0)). [\rho_l(R)]_0 dR \\
    &=\kappa_l \int_{R \in SO(3)}
        f^\dagger(T_{\tilde R^{-1} R}(x_0)). [\rho_l(R)]_0 dR \\
    &=\kappa_l \int_{R' \in SO(3)}
        f^\dagger(T_{R'}(x_0)). [\rho_l(\tilde{R}R')]_0 \tilde{R}^{-1}dR' \\
    &=\kappa_l \int_{R' \in SO(3)}
        f^\dagger(T_{R'}(x_0)). [\rho_l(\tilde{R}).\rho_l(R')]_0 \tilde{R}^{-1}dR'
         \qquad \text{($\rho_l$ is representation)}\\
    &=\kappa_l \rho_l(\tilde{R}).
        \int_{R' \in SO(3)}
            f^\dagger(T_{R'}(x_0)).[\rho_l(R')]_0 \tilde{R}^{-1} dR'
     \qquad \text{(vector interpretation of matrix multiplication (*))}\\
    &=\kappa_l \rho_l(\tilde{R}).
        \int_{R' \in SO(3)}
            f^\dagger(T_{R'}(x_0)).[\rho_l(R')]_0 dR'
     \qquad \text{($\tilde{R}^{-1}$ isometry)} \\
     &= \rho_l(\tilde{R}) .  \mathcal{F}(f)_{l}
\end{align*}

We note that line $(*)$ also proves the equivariance property of the spherical harmonics.

The equivariance of the matrix Fourier transform is proved similarly, and we prove it in the general case of homogeneous spaces in the next section.

\subsection{Bispectrum: Definitions}

The two definitions of Fourier transform (vector-valued and matrix-valued) give rise to two definitions of bispectrum. We focus on the vector-valued case, as the matrix-valued case is tackled in all generality in the next section.

\textbf{The (vector) bispectrum} of a spherical function $f$ is a vector $\beta(f)_{l_1l_2}$ formed of the following scalar entries:
\begin{equation}
    \beta(f)_{l_1l_2}[l]:=\left[\mathcal{F}_0(f)_{l_1} \otimes \mathcal{F}_0(f)_{l_2} \right] C_{l_1l_2} \hat{\mathcal{F}_0(f)_l}^{\dagger} \qquad \in \mathbb{C}; \quad|l_1-l_2| \leq l \leq l_1+l_2.
\end{equation}
where $C_{l_1l_2}$ is the Clebsch-Gordan matrix associated with $\rho_{l_1} \otimes \rho_{l_2}$, and $\hat{\mathcal{F}}_0(f)_l$ is the $(2l_1+1)(2l_2+1)$ vector padded with zeros: $\hat{\mathcal{F}}_0(f)_l:=\left[0, \ldots, 0, \mathcal{F}_0(f)_l, 0, \ldots, 0\right]$, where the number of zeros preceding $\mathcal{F}_0(f)_l$ is $n_p$, and the number succeeding $\mathcal{F}_0(f)_l$ is $n_s$, where $n_p:=\sum_{q=|\ell_1-\ell_2|}^{l-1}(2 q+1)$ and $n_s:=\sum_{q=l+1}^{\ell_1+\ell_2}(2 q+1)$.

Likewise, $\mathcal{F}_0(f)_{l_1} \otimes \mathcal{F}_0(f)_{l_2}$ is a $(2l_1+1)(2l_2+1)$ vector, such that the resulting $\beta(f)_{l_1l_2}[l]$ is indeed a scalar.

\textbf{The (matrix) bispectrum} of a spherical function $f$ is:
\begin{equation}
    \beta(f)_{l_1l_2} = [\mathcal{F}(f)_{l_1} \otimes \mathcal{F}(f)_{l_2}] C_{l_1,l_2} \Big[ \bigoplus_{l \in l_1 \otimes l_2} \mathcal{F}(f)_{l}^{\dagger} \Big] C^{\dagger}_{l_1,l_2} \qquad \in \mathbb{C}^{(2l_1+1)(2l_2 + 1) \times (2l_1+1)(2l_2 +1)}
\end{equation}

% \textbf{The bispectrum for spherical functions} is built from the $(2 l_1+1)(2l_2+1)$-dimensional tensor product vectors $\widehat{\boldsymbol{f}}_{l_1} \otimes \widehat{\boldsymbol{f}}_{l_2}$, which transform according to
% $$
% \widehat{\boldsymbol{f}}_{l_1} \otimes \widehat{\boldsymbol{f}}_{l_2} \mapsto\left(D^{\left(l_1\right)}(R) \otimes D^{\left(l_2\right)}(R)\right) \cdot\left(\widehat{\boldsymbol{f}}_{l_1} \otimes \widehat{\boldsymbol{f}}_{l_2}\right) .
% $$

% \textbf{The non-zero elements of the bispectrum}  for spherical functions are (something like):

% \begin{aligned}
% p_{l_1, l_2}
%     &=\widehat{f}_{l_1} \otimes \widehat{f}_{l_2} \bigoplus_{l \in l_1 \otimes l_2}\widehat{g}_{l_1, l_2, l}^{\dagger} \cdot \widehat{f}_l \qquad \in \mathbb{C}^{D_1 D_2}\\
%     & = \bigoplus_{l \in l_1 \otimes l_2} \sum_{m=-l}^l \sum_{m_1=-l_1}^{l_1} C_{m_1, m-m_1, m}^{l_1, l_2, l} \widehat{f}_{l_1, m_1}^* \widehat{f}_{l_2, m-m_1}^* \widehat{f}_{l_{l, m}} .
% \end{aligned}


% Up to unitary transformation, these invariants are equivalent to the non-vanishing matrix elements of the abstract bispectrum (as already derived in [3] and [1]). Any kernel built from the bispectrum using (प3) as features will be invariant to translation and rotation.

\subsection{Bispectrum: Invariances}

The two definitions of bispectrum are invariant. We focus on the vector-valued case, as the matrix-valued case is tackled in all generality in the next section. For the vector-valued case, we prove that each of the scalar entries forming the bispectrum vector is equivariant.

Define the transformation of $f$ by $\tilde R$ as $(\tilde{R} \ast f)(x) = f({\tilde R}^{-1}\circ x)$. Under this group action, the coefficient $l$ of the vector bispectrum becomes:
\begin{align*}
    \beta(\tilde{R} \ast f)_{l_1l_2}[l]
    &=\left[
    \mathcal{F}(\tilde{R} \ast f)_{l_1} \otimes \mathcal{F}(\tilde{R} \ast f)_{l_2}
    \right]
    C_{l_1l_2} \mathcal{F}_0(\tilde{R} \ast f)_l^{\dagger} \\
    &=\left[
        (\mathcal{F}(f)_{l_1}.\rho_{l_1}(\tilde{R})) \otimes (\mathcal{F}(f)_{l_2}.\rho_{l_2}(\tilde{R}))
        \right]
        C_{l_1l_2} .\mathcal{F}_0(\tilde{R} \ast f)_l^{\dagger}
        \qquad \text{(equivariance on the right with row-vectors)} \\
    &=\left[
        \mathcal{F}(f)_{l_1} \otimes (\mathcal{F}(f)_{l_2})
      \right] .
        \left[
        \rho_{l_1}(\tilde{R}) \otimes \rho_{l_2}(\tilde{R})
        \right] .
        C_{l_1l_2} .
        \mathcal{F}_0(\tilde{R} \ast f)_l^{\dagger}
        \\
     &=\left[
        \mathcal{F}(f)_{l_1} \otimes (\mathcal{F}(f)_{l_2})
      \right] .
        \left[
        \rho_{l_1}(\tilde{R}) \otimes \rho_{l_2}(\tilde{R})
        \right] .
        C_{l_1l_2} .
        \left[ \mathcal{F}_0(f)_l.\bigoplus_{i=|l_1 - l_2|}^{l_1+l_2} \rho_i(\tilde R)\right]^\dagger
        \qquad \text{(only $i=l$ is non-zero?)}
        \\
     &=\left[
        \mathcal{F}(f)_{l_1} \otimes (\mathcal{F}(f)_{l_2})
      \right] .
        \left[
        \rho_{l_1}(\tilde{R}) \otimes \rho_{l_2}(\tilde{R})
        \right] .
        C_{l_1l_2} .
        \left[ \bigoplus_{i=|l_1 - l_2|}^{l_1+l_2} \rho_i(\tilde R)\right]^\dagger
        . \mathcal{F}_0(f)_l^\dagger
        \\
     &=\left[
        \mathcal{F}(f)_{l_1} \otimes (\mathcal{F}(f)_{l_2})
      \right] .
        \left[
        \rho_{l_1}(\tilde{R}) \otimes \rho_{l_2}(\tilde{R})
        \right] .
        C_{l_1l_2} .
        \left[ \bigoplus_{i=|l_1 - l_2|}^{l_1+l_2} \rho_i(\tilde R)\right]^\dagger
        .C_{l_1l_2}^\dagger
        . C_{l_1l_2}  \mathcal{F}_0(f)_l^\dagger
        \qquad \text{(unitary $C_{l_1l_2}$)}\\
     &=\left[
        \mathcal{F}(f)_{l_1} \otimes (\mathcal{F}(f)_{l_2})
      \right]
        . C_{l_1l_2}  \mathcal{F}_0(f)_l^\dagger
        \qquad \text{(Clebsh-Gordan)}\\
     &= \beta(f)_{l_1l_2}[l]
\end{align*}

\section{\color{purple}{Domain $\Omega = (H,~G,~\ast)$ (homogeneous space)}}

By definition, $H$ is a homogeneous space of $G$ if fixing any $x_0 \in H$, the set $T_g(x_0)$ ranges over the whole of $H$ as $g$ ranges over $G$.

The classical example of a homogeneous space is the unit sphere $S_2$ which was the focus of the previous section. The sphere is a homogeneous space of the three-dimensional rotation group $S O(3)$ : taking the North pole as $x_0$, a suitable rotation can move it to any point $x \in S_2$.



\subsection{Fourier Transform: Definition}

\textbf{The Fourier transform} for functions defined over an homogeneous space $H$ of a group $G$ is:
\begin{equation}
\mathcal{F}(f)_\rho = \widehat{f}(\rho) = \int_{g \in G} f\left(T_g(x_0)\right) \rho(g)dg.
\end{equation}

Kondor comments on this equation as follows. Note that except for the trivial case $H=G$, Fourier transforms on homogeneous spaces are naturally redundant: typically $H$ is a much smaller space than $G$, yet a Fourier transform on $H$ has the same number of components as a Fourier transform on the entire group, i.e., as many as the number of irreducible representations of the group. One manifestation of this fact is that we might find that some Fourier components are rank deficient no matter what $f: H \rightarrow \mathbb{C}$ we choose.

Kondor continues. While this destroys Kakarala's uniqueness result, in practice we often find that the bispectrum still furnishes a remarkably rich invariant representation of $f$. We remark that that invariance to right-translation $f^{(z)}(x)=f\left(x z^{-1}\right)$ would be a different matter: there is a variant of the bispectrum which retains the uniqueness property in this case.

Consider the function $\tilde{f}$ defined on the group $G$, such that: $\tilde{f}(g) = f(\pi(g)) = f(T_g.x_0)$, i.e. when $\pi$ is the projection to the quotient space, defined as:
\begin{align*}
    \pi : G \rightarrow H \, \qquad g \rightarrow \pi(g) = T_g.x_0.
\end{align*}

The Fourier transform, for function over groups, gives:
\begin{equation}
\mathcal{F}(\tilde{f})_\rho
=
\widehat{\tilde{f}}(\rho)
=
\int_{g \in G} \tilde{f}(g) \rho(g)dg
=
\int_{g \in G} f(\pi(g)) \rho(g)dg
=
\int_{g \in G} f(T_g.x_0) \rho(g)dg ,
\end{equation}
which is the definition we use for homogeneous space.

Thus, the Fourier transform of a function $f$ defined on a $H$ is the same as the (group) Fourier transform for a lifted function $\tilde f$ defined on the corresponding group.

\subsection{Fourier Transform: Equivariance}

Define the transformation of $f$ by $\tilde g$ as $(\tilde{g} \ast f)(x) = f(T_{\tilde g}^{-1}.x)$. The Fourier transform undergoes a transformation by the same $\tilde g$, as follows, which shows its equivariance.

\vspace{-0.7cm}
\begin{align*}
  \mathcal{F}(\tilde{g} \ast f)_\rho
  &=  \int_{g \in G} (\tilde{g} \ast f)(T_g(x_0))\rho(g) dg \\
  &=  \int_{g \in G} f(T_{\tilde g}^{-1}.T_g(x_0))\rho(g) dg \\
  &=  \int_{g \in G} f(T_{\tilde g^{-1}\circ g}(x_0))\rho(g) dg \\
  &=  \int_{g' \in G} f(T_{g'}(x_0))\rho(\tilde{g}\circ g') dg' \\
  &=  \int_{g' \in G} f(T_{g'}(x_0))\rho(\tilde{g}).\rho(g') dg' \\
  &=  \rho(\tilde{g}) . \int_{g' \in G} f(T_{g'}(x_0))\rho(g') dg' \\
  &=  \rho(\tilde{g}) .\mathcal{F}(f)_\rho
\end{align*}
\vspace{-0.7cm}

Importantly, this equivariance is the same equivariance as the one that emerged for the Fourier transform of function defined over a group.

\subsection{Triple Correlation: Definition}

The triple correlation of a function $f$ defined on $H$ homogeneous for a compact group $G$ is:
\begin{equation}
    T(f)_{g_1, g_2}
        =
        \int_{g \in G}
        f^*(T_g(x_0))
        f\left(T_{g g_1}(x_0)\right) f\left(T_{g g_2}(x_0)\right)
        dg .
\end{equation}
We can rewrite this definition in terms of the function $\tilde f$ that is the lift of function $f$ on the group $G$: $\tilde{f}(g) = f(T_g.x_0) = f(\pi(g))$.
\begin{equation}
    T(f)_{g_1, g_2}
        =
        \int_{g \in G}
        \tilde{f}^*(g)
        \tilde{f}(gg_1) \tilde{f}(gg_2)
        dg ,
\end{equation}
where we recognize the definition of triple correlation on groups applied to function $\tilde{f}$.

\subsection{Triple Correlation: Invariance}
Note that $T(f)$ is a function defined on $H \times H$, and that, like the autocorrelation, the triple correlation is invariant under left translation: $T(f_1)=T(f_2)$ if $f_1(g)=f_2(h g)$.

\begin{proof}
Consider two real signals $f_1, f_2$ defined on a domain $\domain$ equipped with the action of a group $G$, such that $f_2 = T_h[f_1]$ in the sense defined in the previous section.

We show that this implies that $T_h[f_1] = T_h[f_2] $. Taking $g_1, g_2 \in G$, we have:
\begin{align*}
    T(f_2)_{g_1, g_2}
        &= \int_{g \in G}
            f_2^*(T_g(x_0)) f_2\left(T_{g g_1}(x_0)\right) f_2\left(T_{g g_2}(x_0)\right) dg  \\
        &= \int_{g \in G}
            T_h[f_1]^*(T_g(x_0)) T_h[f_1]\left(T_{g g_1}(x_0)\right) T_h[f_1]\left(T_{g g_2}(x_0)\right) dg  \\
        &= \int_{g \in G}
            f_1^*(T_h \circ T_g(x_0)) f_1\left(T_h \circ T_{g g_1}(x_0)\right) f_1\left(T_h \circ T_{g g_2}(x_0)\right) dg  \\
        &= \int_{g \in G}
            f_1^*(T_{hg}(x_0)) f_1\left(T_{hg g_1}(x_0)\right) f_1\left(T_{hg g_2}(x_0)\right) dg  \\
        &= \int_{g \in G}
            f_1^*(T_{g}(x_0)) f_1\left(T_{g g_1}(x_0)\right) f_1\left(T_{g g_2}(x_0)\right) dg \\
        &= T(f_1)_{g_1, g_2}
\end{align*}
where we use the change of variable $hg \rightarrow g$.
\end{proof}

\subsection{Bispectrum: Definitions}

The bispectrum can be defined more generally for signals defined over a space that is homogeneous for a compact group \cite{kakarala2012bispectrum}.

\textbf{The bispectrum for homogeneous spaces of commutative groups} is defined over a pair of irreducible representations $\rho_1, \rho_2$:

\begin{equation}
    \beta (f)_{\rho_1, \rho_2} = \hat{f}_{\rho_1}\hat{f}_{\rho_2}\hat{f}_{\rho_1 \rho_2}^{\dagger}.
    \label{eqn:group-bs}
\end{equation}
The irreducible representations (irreps) of commutative groups map to scalar values.

\textbf{The bispectrum for homogeneous spaces of non-commutative groups} is defined over a pair of irreducible representations $\rho_1, \rho_2$:

\begin{align*}
   \beta (f)_{\rho_1, \rho_2}
   &= [\mathcal{F}(f)_{\rho_1} \otimes \mathcal{F}(f)_{\rho_2}] C_{\rho_1,\rho_2} \Big[ \bigoplus_{\rho \in \rho_1 \otimes \rho_2} \mathcal{F}(f)_{\rho}^{\dagger} \Big] C^{\dagger}_{\rho_1,\rho_2}\\
   &= [\hat{f}{\rho_1} \otimes \hat{f}{\rho_2}] C_{\rho_1,\rho_2} \Big[ \bigoplus_{\rho \in \rho_1 \otimes \rho_2} \hat{f}_{\rho}^{\dagger} \Big] C^{\dagger}_{\rho_1,\rho_2}
    \label{eqn:non-commutative-bs}
\end{align*}

We note that the bispectrum of functions defined over an homogeneous space $H$ has a definition identical to the case of functions defined over a group. The only difference is that the Fourier $\mathcal{F}$ transform used is the Fourier transform for function defined over $H$, and not $G$.

\begin{proof}
By definition, the bispectrum is the Fourier transform of the triple correlation. We compute this Fourier transform and show that we obtain the expression given above.


\begin{align*}
   \beta (f)_{\rho_1, \rho_2}
   &=
   \int_{g_1 \in G} \int_{g_2 \in G}
        T(f)_{g_1, g_2}
            \left[
                \rho_1(g_1) \otimes \rho_2(g_2)
            \right]
            dg_1 dg_2 \\
   &=
    \int_{g_1 \in G} \int_{g_2 \in G}
        \int_{g \in G}
            f^*(T_g(x_0))
            f\left(T_{g g_1}(x_0)\right) f\left(T_{g g_2}(x_0)\right)
            dg
              \left[
                \rho_1(g_1) \otimes \rho_2(g_2)
              \right]
            dg_1 dg_2 \\
\end{align*}


Below, we give the main elements of the proof that shows that computing this Fourier transform indeed yields the same formula as in the case of compact groups.

Note that the group $G$ is often ``bigger'' than the homogeneous space $H$, e.g. $SO(3)$ has dimension 3 while $S^2$ has dimension 2. Any homogeneous space $H$ can be seen as $H = G / G_0$ where $G_0$ is one stabilizer group, e.g. $S^2 = SO(3) / SO(2)$ with $G_0 = SO(2)$.

A function $f$ defined on $H = G / G_0$ can be lifted to a function defined on $G$, as $\tilde{f}(g) = f(\pi(g))$ where $\pi$ is the canonical projection associated with the quotient space structure. The lifted function $\tilde{f}$ is invariant under translations by elements of the stabilizer $G_0$:
\begin{equation}
    \tilde{f}(g_0.g) = \tilde{f}(g), \qquad \text{for all $g_0 \in G_0$},
\end{equation}
This invariance induces a property on its fourier transform:
\begin{equation}
    \mathcal{F}(\tilde{f})(\rho) = \mathcal{F}(\tilde{f})(\rho) \rho(g_0), \qquad \text{for all $g_0 \in G_0$.}
\end{equation}

By integration over $G_0$, we get:
\begin{equation}
    \mathcal{F}(\tilde{f})(\rho) =  \mathcal{F}(f)(\rho).P_0(\rho)\quad \text{where: } P_0(\rho) = \int_{g_0 \in G_0} \rho(g_0),
\end{equation}
where $P_0$ is a projection operator: $P_0 . P_0 = P_0$.

Next, one can prove that see that tensor products of projection operators decompose according to the Clebsh-Gordon formula:
$$
\begin{gathered}
P_0(\sigma) \otimes P_0(\delta)= \\
C_{\sigma \delta}\left[P_0\left(\omega_1\right) \oplus \cdots \oplus P_0\left(\omega_k\right)\right] C_{\sigma \delta}^{\dagger}\left[P_0(\sigma) \otimes P_0(\delta)\right] .
\end{gathered}
$$
This ensures that the bispectrum formula on compact groups is valid for functions whose domain is a homogeneous space as well as those defined on groups.
\end{proof}

\subsection{Bispectrum: Invariance}

The proof of invariance of the bispectrum for function defined over (commutative and non-commutative) domains does not use the expression of the Fourier transform on groups. Rather, it only uses the fact that the Fourier transform on groups is equivariant, with equivariance factor $\rho(\tilde{g})$. Since the Fourier transform on homogeneous spaces is equivariant, with the same equivariance factor, and since the bispectra definition is the same, we get the invariance of the bispectrum with the same derivation.

We note that the proofs of bispectrum invariance of the previous section does not depend on the specific fact that we are using the group of rotations $SO(3)$. Therefore, they can be readily generalized to create a more general vector-valued bispectrum on homogeneous spaces. The remaining TODO is to define the vector-valued Fourier transform in the general case, and prove its equivariance.

\section{\color{deepred}{Domain $\Omega = (S^1 \times \mathbb{R}_+^*,~SO(2),~\ast)$ (disk rotations)}}


\subsection{Fourier Transform: Definition}

Functions defined over a disk $D$ can be expanded on the disk harmonics as:
\begin{equation}
    f(r, \theta)=\sum_{m, n} \hat{f}(m, n) d_{n m}(r, \theta),
\end{equation}
where $ d_{m m}(r, \theta)$ is the disk harmonics of coefficients $m, n$ given by:
\begin{equation}
    d_{nm}(r, \theta)= J_m(2 \pi l_{n m} r) \exp (im \theta)
\end{equation}

\textbf{The (vector) Fourier transform} $\hat{f}(m, n)$ for functions defined over a disk $D$ is:
\begin{align*}
  \mathcal{F}(f)_{nm}
    &= \frac{1}{a_{n m}} \int_0^{2 \pi} \int_0^1 f(r, \theta) J_m(2 \pi l_{n m} r) \exp(-i m \theta) r d r d\theta \\
    &=\left\{\begin{array}{c}
\frac{1}{2 \pi i m} a_{n m} \int_0^{2 \pi} F(l_{n m}, \phi) \exp(-i m \phi) d \phi \quad \text{for $l\neq 0$} \\
F(0,0) / \pi \quad \text{for $l = 0$}
\end{array}\right\},
\end{align*}

where $F(l_{n m}, \phi) $ is the traditional Fourier transform of $f(r, \theta)$ written in polar coordinates.

\subsection{Fourier Transform: Equivariance}

Define the transformation of $f$ by $\tilde \theta$ as $(\tilde{\theta} \ast f)(r, \theta) = f(r, \tilde \theta +\theta )$. The Fourier transform undergoes a transformation by the same $\tilde \theta$, as follows, which shows its equivariance.

\begin{align*}
    \mathcal{F}(\tilde{\theta} \ast f)_{nm}
    &= \frac{1}{a_{n m}}
        \int_0^{2 \pi} \int_0^1
            (\tilde{\theta} \ast f)(r, \theta) J_m(2 \pi l_{n m} r) \exp(-i m \theta)
                r dr d\theta \\
    &= \frac{1}{a_{n m}}
        \int_0^{2 \pi} \int_0^1
            f(r, \tilde \theta + \theta) J_m(2 \pi l_{n m} r) \exp(-i m \theta)
                r dr d\theta \\
    &= \frac{1}{a_{n m}}
         \int_{\tilde{\theta}}^{2 \pi+\tilde{\theta}} \int_0^1
            f(r, \theta') J_m(2 \pi l_{n m} r) \exp(-i m (\theta' - \tilde{\theta}))
                r dr d\theta'\\
    &= \frac{1}{a_{n m}}
         \int_{\tilde{\theta}}^{2 \pi+\tilde{\theta}} \int_0^1
            f(r, \theta') J_m(2 \pi l_{n m} r) \exp(-i m \theta')\exp(+im\tilde{\theta})
                r dr d\theta' \\
    &= \exp(im\tilde{\theta}) .\frac{1}{a_{n m}}
         \int_{\tilde{\theta}}^{2 \pi+\tilde{\theta}} \int_0^1
            f(r, \theta') J_m(2 \pi l_{n m} r) \exp(-i m \theta')
                r dr d\theta' \\
    &= \exp(im\tilde{\theta}) .\frac{1}{a_{n m}}
         \int_{0}^{2 \pi} \int_0^1
            f(r, \theta') J_m(2 \pi l_{n m} r) \exp(-i m \theta')
                r dr d\theta' \\
    &= \exp(im\tilde{\theta}) . \mathcal{F}(f)_{nm}
\end{align*}


\subsection{Bispectrum: Definition}

We generalize the bispectrum for signals defined over a space that is the direct product of $H$, homogeneous for a compact group, and $R$ invariant for that group. We do it here in the case of the disk.

\textbf{The bispectrum for the disk} is defined as:
\begin{equation}
    \beta(f)_{nmm'} = \mathcal{F}(f)_{nm}.\mathcal{F}(f)_{nm'}.\mathcal{F}(f)_{n(m+m')}^\dagger,
\end{equation}
where we observe that the third term only depends on the indices $m$, $m'$ that correspond to the homogeneous part of the disk, i.e. the one that transform with rotations in $SO(2)$. We recognize that we add these indices, since $SO(2)$ is a commutative group.

\subsection{Bispectrum: Invariance}


Define the transformation of $f$ by $\tilde \theta$ as $(\tilde{\theta} \ast f)(r, \theta) = f(r, \tilde \theta +\theta )$. Under this group action, the bispectrum becomes:
\begin{align*}
     \beta(\tilde{\theta} \ast f)_{nmm'}
        &= \mathcal{F}(\tilde{\theta} \ast f)_{nm}.\mathcal{F}(\tilde{\theta} \ast f)_{nm'}.\mathcal{F}(\tilde{\theta} \ast f)_{n(m+m')}^\dagger\\
        &= \exp(im\tilde{\theta}) \mathcal{F}(f)_{nm}.\exp(im'\tilde{\theta}) \mathcal{F}(f)_{nm'}.\exp(-i(m+m')\tilde{\theta}) \mathcal{F}(f)_{n(m+m')}^\dagger\\
        &=\beta(f)_{nmm'},
\end{align*}
which proves its invariance.

\section{\color{deepred}{Domain $\Omega = (S^2 \times \mathbb{R}_+^*,~SO(3),~\ast)$ (ball rotations)}}

\subsection{Fourier Transform: Definition}

Functions defined over a ball $B$ can be expanded on the surface spherical harmonics as:
\begin{align*}
       f(r, \theta, \phi)
        &=\sum_{\ell=0}^{\infty} \sum_{m=-\ell}^{\ell}
            \hat{f}_{\ell}^m r^{\ell} Y_{\ell}^m(\theta, \phi)
\end{align*}


\textbf{The (vector) Fourier transform} $\hat{f}_\ell$ for functions defined over a ball $B$ is (check?):
\begin{align*}
  \mathcal{F}(f)_{l}
    &\propto \int_0^{2 \pi}\int_0^{\pi} \int_0^1 f(r, \theta, \phi)^\dagger  \left( r^{\ell} Y_{\ell}(\theta, \phi)\right) r d r \cos(\theta) d\theta d\phi \\
    &= \sqrt{\frac{2l+1}{4\pi}}  \int_0^{2 \pi}\int_0^{\pi} \int_0^1 f(r, \theta, \phi)^\dagger  \left( r^{\ell}  [\rho_l(\phi, \theta, \gamma)]_0 \right) r d r \cos(\theta) d\theta d\phi \\
    &= \sqrt{\frac{2l+1}{4\pi}}  \int_0^{2 \pi}\int_0^{\pi} \int_0^1 f(r, T_R(x_0))^\dagger  \left( r^{\ell}  [\rho_l(\phi, \theta, \gamma)]_0 \right) r d r \cos(\theta) d\theta d\phi \\
    &=\sqrt{\frac{2l+1}{4\pi}} \int_0^1 \left(\int_{R \in SO(3)} f(r, T_R(x_0))^\dagger [\rho_l(R)]_0. dR\right) r^{l+1} dr
\end{align*}

\subsection{Fourier Transform: Equivariance}

Define the transformation of $f$ by $\tilde R$ as $(\tilde{R} \ast f)(r, \theta, \phi) = f(r, \tilde R (\theta, \phi)) = f(T_{\tilde R}^{-1}.u)$. The (vector) Fourier transform undergoes a transformation by the same $\tilde R$, as follows, which shows its equivariance.

\begin{align*}
    \mathcal{F}(\tilde{R} \ast f)_{l}
    &=\kappa_l \int_0^1 \left(\int_{R \in SO(3)} (\tilde{R} \ast f)(r, T_R(x_0))^\dagger [\rho_l(R)]_0. dR\right) r^{l+1} dr\\
    &=\kappa_l \int_0^1 \left(\int_{R \in SO(3)} f(r, T_R^{-1}.T_R(x_0))^\dagger [\rho_l(R)]_0. dR\right) r^{l+1} dr\\
    &=\kappa_l \int_0^1 \left(\int_{R \in SO(3)} f(r, T_{R^{-1}R}(x_0))^\dagger [\rho_l(R)]_0. dR\right) r^{l+1} dr\\
    &=\kappa_l \int_0^1 \left(\int_{R' \in SO(3)} f(r, T_{R'}(x_0))^\dagger [\rho_l(\tilde{R}R')]_0. \tilde{R}^{-1}dR'\right) r^{l+1} dr\\
    &=\kappa_l \int_0^1 \left(\int_{R' \in SO(3)} f(r, T_{R'}(x_0))^\dagger [\rho_l(\tilde{R})\rho_l(R')]_0. \tilde{R}^{-1}dR'\right) r^{l+1} dr\\
    &=\rho_l(\tilde{R}) . \kappa_l \int_0^1 \left(\int_{R' \in SO(3)} f(r, T_{R'}(x_0))^\dagger [\rho_l(R')]_0. \tilde{R}^{-1}dR'\right) r^{l+1} dr\\
    &=\rho_l(\tilde{R}) . \kappa_l \int_0^1 \left(\int_{R' \in SO(3)} f(r, T_{R'}(x_0))^\dagger [\rho_l(R')]_0. dR'\right) r^{l+1} dr\\
    &=\rho_l(\tilde{R}) . \mathcal{F}(f)_{l} \\
\end{align*}


\subsection{Bispectrum: Definitions}

We define the bispectrum as in the spherical function case: the only difference is the type of Fourier transform used, which is the one defined above. Then, both the definition and the proof of invariance of the bispectrum follow.

\textbf{The (vector) bispectrum} of a ball function $f$ is a vector $\beta(f)_{l_1l_2}$ formed of the following scalar entries:
\begin{equation}
    \beta(f)_{l_1l_2}[l]
    :=
    \left[
    \mathcal{F}(f)_{l_1} \otimes \mathcal{F}(f)_{l_2}
    \right]
    C_{l_1l_2} \mathcal{F}_0(f)_l^{\dagger}
    \qquad \in \mathbb{C}; \quad|l_1-l_2| \leq l \leq l_1+l_2.
\end{equation}
where $C_{l_1l_2}$ is the Clebsch-Gordan matrix associated with $\rho_{l_1} \otimes \rho_{l_2}$, and $\mathcal{F}_0(f)_l$ is the $(2l_1+1)(2l_2+1)$ vector padded with zeros: $\mathcal{F}_0(f)_l:=\left[0, \ldots, 0, \mathcal{F}_l, 0, \ldots, 0\right]$, where the number of zeros preceding $\mathcal{F}_i$ is $n_p$, and the number succeeding $\mathcal{F}_i$ is $n_s$, where $n_p:=\sum_{q=|k-\ell|}^{i-1}(2 q+1)$ and $n_s:=\sum_{q=i+1}^{k+\ell}(2 q+1)$.

\subsection{Bispectrum: Invariance}

The proof of invariance is identical to the proof of invariance for the $S^2$ case, since the Fourier transform on the ball has the same equivariance property as the Fourier transform on the sphere.

\section{\color{deepred}{Domain $\Omega = (H \times R,~G,~\ast)$ (homogeneous $\times$ invariant space)}}

\subsection{Fourier Transform: Definition}


\textbf{The (matrix) Fourier transform} for functions defined over $H \times R$ is:
\begin{align*}
  \mathcal{F}(f)_{\rho}
    &= \int_{r \in R} \int_{g \in G} f\left(T_g(x_0)\right) \rho(g)dg. dr \\
    &= \int_{r \in R} \int_{g \in G} f\left(r, T_g(x_0)\right) \rho(g)dg. dr
\end{align*}

TODO: find how integral on $r$ should depends on $\rho$, by comparing to the disk and ball cases, in order to be a valid Fourier transform.

\textbf{The (vector) Fourier transform} might exist in this case too, by analogy with the disk and ball cases.

\subsection{Fourier Transform: Equivariance}

Define the transformation of $f$ by $\tilde g$ as $(\tilde{g} \ast f)(r, x) = f(r, \tilde g^{-1} \ast x) = f(T_{\tilde g}^{-1}.x)$. The (vector) Fourier transform undergoes a transformation by the same $\tilde R$, as follows, which shows its equivariance.

\begin{align*}
  \mathcal{F}(\tilde{g} \ast f)_{\rho}
    &= \int_{r \in R} \int_{g \in G} (\tilde{g} \ast f)\left(T_g(x_0)\right) \rho(g)dg. dr \\
    &= \int_{r \in R} \int_{g \in G} (\tilde{g} \ast f)\left(r, T_g(x_0)\right) \rho(g)dg. dr \\
    &= \int_{r \in R} \int_{g \in G} f(r, T_{\tilde{g}}^-1.T_g(x_0)) \rho(g)dg. dr \\
    &= \int_{r \in R} \int_{g \in G} f(r, T_{\tilde{g}^-1g}(x_0)) \rho(g)dg. dr \\
    &= \int_{r \in R} \int_{g' \in G} f(r, T_{g'}(x_0)) \rho(\tilde{g}g')dg'. dr \\
    &= \int_{r \in R} \int_{g' \in G} f(r, T_{g'}(x_0)) \rho(\tilde{g})\rho(g')dg'. dr \\
    &= \rho(\tilde{g}). \int_{r \in R} \int_{g' \in G} f(r, T_{g'}(x_0)) \rho(g')dg'. dr \\
    &= \rho(\tilde{g}) \mathcal{F}(f)_{\rho}
\end{align*}


Note that the equivariance does not depend on the dependency of the integral on $\rho$, thus any ``variant'' of the Fourier transform above will be equivariant.

\subsection{Bispectrum: Definition}

The definition of the bispectrum is identical to the its definition in the homogeneous case: only the Fourier transform used varies.

\subsection{Bispectrum: Invariance}

The invariance of the bispectrum follows from the proof of invariance on homogeneous spaces, since the Fourier transforms on $H \times R$ and on $H$ have the same equivariance properties.

\section{\color{deepred}{Domain $\Omega = (\mathbb{Z}^2, G, $ with Steerable CNN}}

Consider the semi-direct product group $G = \mathbb{Z}^2 \ltimes H$, where $\mathbb{Z}_2$ is the domain. We consider a space of functions $f:\mathbb{Z}^2 \mapsto \mathbb{R}^K$ equipped with the representation $T$ of $G$ induced by a representation $\rho_H$ of $H$ on the fiber $\mathbb{R}^K$, and is denoted $T=\operatorname{Ind}_H^G \rho$. For a function $f$, the action of an element $(t, r) \in \mathbb{Z}^2 \ltimes H$ is:
\begin{equation}
    \left[T_{t r} f\right](x)=\rho_H(r)\left[f\left((t r)^{-1} x\right)\right] \in \mathbb{R}^K
\end{equation}
Even if this uses notions of representations (e.g. $\rho_H$ and $\pi$, these are just there to explicit the action on the signal).

\subsection{Fourier Transform: Definition}

\textbf{The Fourier transform} for functions defined over a domain $\Omega$ but evolving according to the induced representation is defined as:
\begin{equation}
\mathcal{F}(f)_\rho = \widehat{f}(\rho ) = \int_{g \in G} \left[T_{t r} f\right](x_0) \otimes \rho(rt)^\dagger dg.
\end{equation}

Here, we use the notation $\rho$ to denote the irreducible representation at which the Fourier coefficient is computed. It is NOT the induced representation of the full group $G = Z_2 \times H$, usually denoted $\pi$ in the steerable CNNs framework, nor $\rho_H$ the representation of the group $H$.

\begin{remark}
We have: $\left[T_{t r} f\right](x_0) \in \mathbb{R}^K$ and $\pi(rt)$ can be a matrix of any dimension, as it depends on the dimension of the irreducible representation.

Thus, what does the product $ \left[T_{t r} f\right](x_0)\pi(rt)$ even mean? It asks the question of what is the Fourier transform of a vector-valued signal. We add a tensor-product. Its intuition is that, at each $g \in G$ we build the $K$-vector that has a matrix at its $k$-component: specifically, $\left(\left[T_{t r} f\right](x_0)\right)_k$ the $k-th$ scalar component of the transformed signal that multiplies the $D \times D$ matrices $\rho(g)$

This is NOT the case where the Fourier transform of a vector-valued signal is defined as the concatenation of the Fourier transform of each of its scalar-valued components. Because, we would not know what group action to use on a signal scalar component $f_k$ of $f$. (unless the restriction of the action on $f$ to only its component $k$ is also an action, which I doubt).
\end{remark}


\subsection{Fourier Transform: Equivariance}


\subsection{Steerable Bispectrum: After G-Convolution on Intertwiners Basis}

We consider a case where the output of the G-convolution has the same domain as the input. Then we can use the same basis functions.

Consider a group $G = \mathbb{Z}^2 \times H$ and $Y$ a vector of stacked functions that form a steerable basis of $\mathbb{Z}^2$ for the group $G$. We need the basis to be steerable for the whole group, because the whole group acts on the domain of the signal.

\begin{remark}[Which basis of functions to decompose the filter?]
    In Steerable CNNs, the filter $\Psi$ is decomposed on a basis of functions $\Psi_l$. Yet, this might not be a basis of steerable functions. It is a basis of functions for the vector space $\text{Hom}(\pi, \rho)$ where $\pi$ is the regular representation (action on the signal's domain) and $\rho$ is a representation of $H$ of our choice. Here, we choose $\rho$ to be the trivial representation for now.
    Additionally the basis of functions should be indexed by the irreducible representations of the group — in order to make sense for the downstream computations of the bispectrum.
\end{remark}

The convolution is written as a spatial convolution in Steerable CNNs, which allows the output's domain to match the input's domain
\begin{align*}
    \Theta(x_0) = (\Psi \ast f) (x_0) = \Psi^T. T_{x_0}f
\end{align*}
where $\Psi = [\Psi_{11}, ..., \Psi_{ss}]$ and $f = [f_{11}, ..., f_{ss}]$ are the values of the filter/signal in spatial format flattened into vectors, while $T_{x_0}f$ is the vector of values of a translated version of the signal.

We write the convolution using the sum over the domain. In the implementation of Steerable CNNs, the signal is translated by $x_0$, not the filter: we apply the filter bank to translated versions of the signal. For the sake of the argument, we perform a change of variable to extract the (possibly steerable??) basis out of the expression of $\Psi$.
\begin{align*}
    \Theta(x_0) = (\Psi \ast f) (x_0)
    &= \sum_{x \in \Omega_s} \Psi(x).f(x_0 - x)\\
    &= \sum_{x' \in \Omega_s^{+x_0}}
        \Psi(x_0 - x').f(x'),
        \quad \text{(note: the domain changes through the change of variable)}\\
    &= \sum_{x' \in \Omega_s^{+x_0}}
        \Psi(x_0 - x').f(x'),\\
    &= \sum_{x' \in \Omega_s^{+x_0}}
        \left(\sum_l \hat{w}^{(l)} Y_l(x_0 - x')\right).f(x'),
        \quad\text{(decomposition on the intertwiner's basis, denoted $Y$)}\\
    &= \sum_{x' \in \Omega_s^{+x_0}}
        \left(\hat{w}^T. Y(x_0 - x')\right).f(x'),
        \quad\text{(inner-product vectorial notation)}\\
    &= \sum_{x' \in \Omega_s^{+x_0}}
        \left(\hat{w}^T. \rho^Y(-x')Y(x_0)\right).f(x'),
        \quad\text{([!] assume the intertwiner's basis is steerable, by $\rho^Y$)}\\
    &= \sum_{x' \in \Omega_s^{+x_0}}
        \left(\hat{w}^T. \rho^Y(-x')Y(x_0)\right).f(x')\\
    &= \left(\sum_{x' \in \Omega_s^{+x_0}}
        \hat{w}^T. \rho^Y(-x').f(x')\right) Y(x_0) \in \mathbb{C}\\
    &= \hat{\Theta}^T.Y(x_0),\quad\text{(where $\hat{\Theta}, Y(x_0) \in \mathbb{C}^L$)} ,\\
    &=\sum_{l=1}^L \hat{\Theta}_l.Y_l(x_0)
\end{align*}
such that: $\hat{\Theta}_l = \left(\sum_{x' \in \Omega_s^{+x_0}} \hat{w}^T. \rho^Y(-x').f(x')\right)_l$, i.e. the $l$-th coordinate of this row vector.

This would work if we can prove: (i) that the intertwiner basis is steerable for $G$ (see [!]), and (ii) that it is indexed by the irreducible representations. For [!],  if the itertwiners' basis are not steerable, maybe we could express them in a steerable basis.

We start with (i). We consider the trivial representation of $H$ in the definition of the intertwiners. THe vector space of intertwiners thus verifies the equation:
\begin{equation}
    Y = Y.\pi(h),\quad \forall h \in H,
\end{equation}
where $Y = [Y_{11}, ..., Y_{ss}] \in \mathbb{C}^{1 \times s^2}$ is a intertwiner, written as the set of its values on the domain. $\pi(h) \in \mathbb{C}^{s^2.s^2}$, the matrix that can be applied to the values of a signal $f = [f_{11}, ..., f_{ss}] \in \mathbb{C}^{s^2}$ and amounts to rotating it. It is known as the regular representation.

The equation above means that the solutions $Y$ are the set of eigenvectors of eigenvalue 1 for every matrix of the representation $\pi$: an invariant subspace of the representation $\pi$, i.e. a sub-representation.


\begin{remark}
    We can rewrite this equation inn terms of coordinates on the domain:
\begin{equation}
    Y(x_0) = \sum_{x \in \in \Omega_s} Y(x).\pi(h)_{xx_0},
\end{equation}
by writing out the matrix multiplication explicitly and indexing the values of $Y$ by the pixel coordinate where that value is located. This equation holds for one element $Y$ of the set of basis functions. We need the whole set to prove the steerability.
\end{remark}

Can the basis of the intertwiners be indexed by the irreps of $H$? To answer this question, we count how many basis functions there are. Writing $\rho_0$ is the trivial representation of $H$ on the fibers, and following the computation after Eq (11) in Steerable CNNs 2017, we have:
\begin{equation}
    L = \text{dim}~\text{Hom}^H(\pi, \rho_0) = \sum_{\text{irreps }\psi } m^\pi_{\psi}.m^{\rho_0}_{\psi} = m^\pi_{\phi_0}
\end{equation}
since the multiplicity of the trivial representation $\rho_0$ is 0 except on itself where it is 1. In the above equation, $m^\pi_{\phi_0}$ is the multiplicity of the representation $\pi$ of $H$ on the trivial irreps of the group $H$. This multiplicity is $3$, as metnioned in Table 1 of the paper.

We take the example of Table 1 of the Steerable CNNs paper. $H = D_4$ the roto-reflexion group. We decompose $\pi_0$ the representation of $G$ on the irreducible representations of $H$. Actually, we consider $\pi_0$ as a representation of $H$ and decompose this one. There are 3 elementary basis filters for the $\pi_0$, meaning that the dimension of Hom is 3, which is different from 5, the number of irreps of $D_4$.

Thus, the number of basis functions in $\text{Hom}^H(\pi, \rho_0)$ is not equal to the number of irreps. We would not get the Fourier coefs as we need them for the bispectrum.

Maybe we could define another bispectrum?

\begin{remark}[If $\pi$ had been the regular representation of $H$]
$\pi$ is probably not the regular representation, because the reg representation has a fixed dimension and here it works for any s, it has shape $s^2 \times s^2$.

From Wikipedia, we have: If G is a finite group and K is the complex number field, the regular representation decomposes as a direct sum of irreducible representations, with each irreducible representation appearing in the decomposition with multiplicity its dimension. Thus, the dimension of the intertwining space is 1, which is different from the number of irreps of the group. Thus, we cannot get the spectral coefficients by reading out the coefficients on the intertwining basis.

On the other hand, why does the Steerable CNN paper talk about basis filter on the irreps?! As if there was one basis filter (i.e. a basis function) per irreps.

Answer: this is the set of intertwiners when the $\rho$ chosen is an irreps. Which is our case, since our $\rho_0$ is the trivial representation. Contrarily to what we just proved, the basis for $\rho_0$, which is $A1$ has three elements.
\end{remark}

\subsection{Bispectrum: from the intertwining coefs?}

Goal: write the output of the convolution of a steerable CNN in terms of its “spectral” coefficients $\hat \Theta_l$ on a steerable basis of the domain.
Illustrative example: G = Z2 x D4, i.e. the example in the steerable CNN paper, where D4 is non-commutative.

Remark: In a steerable CNN, the filter is expressed as “spectral” coefficients on a basis.
However, they do NOT use a steerable basis;
instead, they use on a intertwiners’ basis of H=D4 (see top of p4 of the Steerable CNNs paper).
Let’s assume for now that the intertwiners’ basis is a steerable basis.

Result: The output of the steerable convolution can be written in terms of its “spectral” coefficients $\hat \Theta_l$on the intertwiner basis (screenshot & details in overleaf).

Next step: Compute the bispectrum from the spectral coefficients $\hat \Theta_l$ .

If we want to use an existing formula of the bispectrum:
we need the L spectral coeffs $\hat \Theta_l$ to correspond to the M Fourier coeffs of Theta  — where M is the number of irreps of the group.
--> we need, at least, to have L = M.
However, that is not the case for the example of D4: we have L=3 spectral coeffs but M=5 irreps of D4 (see Table 1 in Steerable CNNs).
Thus, we need a new formula for the bispectrum, written in terms of the $\hat \Theta_l$
Questions:
---> How do we generalize the bispectrum, for a non-commutative group H, when we don’t have the Fourier coefs of the signal, but rather its coefs on an intertwining basis? (that is, maybe, also a steerable basis?)
---> If we find a formula for it, what do we expect it to be invariant to? I’d say invariant to action of H on the domain.

Since we would expect this bispectrum to be invariant to the action of H on the domain, we compute the TC for this action, and then take its Fourier transform.

This is the TC for the action of rotation on the domain. which is not a homogeneous domain. Maybe we could take the disk bispectrum then? And write everything in disk harmonics? The disk harmonics are steerable by $\rho(\theta) = \exp(-im\theta)$ just as the ring harmonics.

Let's rewrite the output of the steerable convolution using the disk harmonics.
\begin{align*}
    \Theta(x_0) = (\Psi \ast f) (x_0)
    &= \sum_{x \in \Omega_s} \Psi(x).f(x_0 - x)\\
    &= \sum_{x' \in \Omega_s^{+x_0}}
        \Psi(x_0 - x').f(x'),
        \quad \text{(note: the domain changes through the change of variable)}\\
    &= \sum_{x' \in \Omega_s^{+x_0}}
        \left(\sum_l \hat{w}^{(l)} Y_l(x_0 - x')\right).f(x'),
        \quad\text{(decomposition on the intertwiner's basis, denoted $Y$)}\\
    &= \sum_{x' \in \Omega_s^{+x_0}}
        \left(\sum_l \hat{w}^{(l)} \sum_{m,n}\hat{Y_l}(m, n) d_{m, n}(r_{x_0 - x'}, \theta_{x_0 - x'})\right).f(x'),
        \quad\text{(decomposition on the harmonics basis, denoted $Y$)}\\
    &= \sum_{m,n} \left(\sum_{x' \in \Omega_s^{+x_0}}
        \sum_l \hat{w}^{(l)} \hat{Y_l}(m, n).f(x')\right) d_{m, n}(r_{x_0 - x'}, \theta_{x_0 - x'}),
        \quad\text{(getting the coefficients on the harmonics)}
    % &= \sum_{x' \in \Omega_s^{+x_0}}
    %     \left(\hat{w}^T. Y(x_0 - x')\right).f(x'),
    %     \quad\text{(inner-product vectorial notation)}\\
    % &= \sum_{x' \in \Omega_s^{+x_0}}
    %     \left(\hat{w}^T. \rho^Y(-x')Y(x_0)\right).f(x'),
    %     \quad\text{([!] assume the intertwiner's basis is steerable, by $\rho^Y$)}\\
    % &= \sum_{x' \in \Omega_s^{+x_0}}
    %     \left(\hat{w}^T. \rho^Y(-x')Y(x_0)\right).f(x')\\
    % &= \left(\sum_{x' \in \Omega_s^{+x_0}}
    %     \hat{w}^T. \rho^Y(-x').f(x')\right) Y(x_0) \in \mathbb{C}\\
    % &= \hat{\Theta}^T.Y(x_0),\quad\text{(where $\hat{\Theta}, Y(x_0) \in \mathbb{C}^L$)} ,\\
    % &=\sum_{l=1}^L \hat{\Theta}_l.Y_l(x_0)
\end{align*}
To conclude, we would need the harmonics to be steerable by the translations.



\subsection{Triple Correlation: Definition}

We define a new notion of triple correlation that works for vectorial function eqquiped with the action of a semi-direct group $\mathbb{Z}^2 \ltimes H$ . In that context, the G-convolution can be conveniently computed using the steerable group convolution. After the first layer of a steerable G-convolution, the function is not a function over the group (as is the case for the classic G-conv) but it is still a function over the $\mathbb{Z}^2$ albeit equipped with an induced group action, i.e. an action that does not only act on the domain $\mathbb{Z}^2$ of the function, but also transforms its values.

We wonder whether we can generalize the triple correlation to this case. We show below that we can.

We modify the definition of the triple correlation by inverting the order of $gg_1$ to $g_1g$:


We define a triple correlation applicable to this class of functions:
\begin{equation}
    T(f)_{g_1g_2}
        = \int_{\mathbb{Z}^2}
            \int_H
                T_{rt}[f]^\dagger (x_0) T_{r_1t_1rt}[f](x_0) T_{r_2t_2rt}[f](x_0)
                    drdt \in \mathbb{R}^K,
\end{equation}
where $g = rt$, $g_1 = r_1t_1$ and $g_2 = r_2 t_2$.

We note that the triple correlation is a vector in $\mathbb{R}^K$, i.e. in the space representing the fiber where $K$ is the number of channels. We also note that we only use one element of the domain $x_0 \in \mathbb{Z}^2$ to define the TC: this is a similar approach as what was done for the homogeneous case. The domain is homogeneous for the action of $G$. By varying $g \in G$, we will see every value of the function.

\subsection{Triple Correlation: Invariance}

We prove invariance of the TC for the steerable case.

We begin with a Lemma:

\begin{lemma}
Composing the group actions corresponding to the induced representation:
    \begin{equation}
    T_{rt}[T_{r't'}[f]] (x)
        = \rho(r) \left[[T_{r't'}[f]]\left((t r)^{-1} x\right)\right]
        = \rho(r)  \rho(r') f((t'r')^{-1}(tr)^{-1}x) \in \mathbb{R}^K
\end{equation}
\end{lemma}


Consider a second function that is obtained from $f$ by an action of a group element $h = r_ht_h$, i.e. we have: $f_2 = T_h[f] = T_{r_ht_h}[f]$. I.e. $f_2$ is in the orbit of $f$ for the action of the induced representation.

\begin{proof}

We prove that $T(f_2)_{g_1g_2} = T(f)_{g_1g_2} $, which proves the invariance of the TC for the action of the induced representation.

\begin{align*}
    T(f_2)_{g_1g_2}
        &= \int_{\mathbb{Z}^2} \int_H
            T_{rt}[f_2]^\dagger (x_0)
                T_{r_1t_1rt}[f_2](x_0)
                    T_{r_2t_2rt}[f_2](x_0)
                    dtdr \\
        &= \int_{\mathbb{Z}^2} \int_H
            T_{rt}[T_{r_ht_h}[f]]^\dagger (x_0)
                T_{r_1t_1rt}[T_{r_ht_h}[f]](x_0)
                    T_{r_2t_2rt}[T_{r_ht_h}[f]](x_0)
                    dtdr \\
        &= \int_{\mathbb{Z}^2} \int_H
            \left(\rho(r)\rho(r_h)
            f((t_hr_h)^{-1}(tr)^{-1}x_0)\right)^\dagger \\
                &\qquad\qquad\qquad.\left(
                \rho(r_1)\rho(r)\rho(r_h)
                f((t_hr_h)^{-1}(tr)^{-1}(t_1r_1)^{-1}x_0)\right)\\
                &\qquad\qquad\qquad.\left(
                \rho(r_2)\rho(r)\rho(r_h)
                f((t_hr_h)^{-1}(tr)^{-1}(t_2r_2)^{-1}x_0)\right) dtdr
\end{align*}
We will prove the rest using the example SE(n) of a semi-direct product group $\mathbb{Z}_2 \ltimes H$.

We recall that, e.g. in SE(n):
\begin{itemize}
    \item the inversion writes: $(tr)^{-1} = (r^{-1}, R^{-1}(-t))$
    \item and the action on an element of the domain writes: $(R, t)x = Rx+t$ .
\end{itemize}
Note: for a general subgroup $H$, replace R by $\phi(r)$ i.e. a representation of $H$ on $\mathbb{Z}_2$ that transforms $r \in H$ into something that acts on $\mathbb{Z}^2$.

Note2: Triple check this.

We expand this:
\begin{align*}
    (t_hr_h)^{-1}(tr)^{-1} x_0
        &= R_h^{-1}(R^{-1}x_0 - R^{-1}t) - R_h^{-1}t_h \\
        &= R_h^{-1}R^{-1}x_0 - R_h^{-1}R^{-1}t- R_h^{-1}t_h
\end{align*}
and:
\begin{align*}
     (t_hr_h)^{-1}(tr)^{-1}(t_1r_1)^{-1} x_0
     &= (t_hr_h)^{-1}(tr)^{-1} (R_1^{-1}x_0 - R_1^{-1}t_1) \\
     &= R_h^{-1}(R^{-1}R_1^{-1}x_0 - R^{-1}R_1^{-1}t_1 - R^{-1}t) - R_h^{-1}t_h \\
     &= R_h^{-1}R^{-1}R_1^{-1}x_0 - R_h^{-1}R^{-1}R_1^{-1}t_1 - R_h^{-1}R^{-1}t - R_h^{-1}t_h \\
\end{align*}

We first perform the change of variable $r' = rr_h$, and $t'=Rt_h+t$  which gives:
\begin{align*}
    \rho(r)\rho(r_h) &= \rho(r')
\end{align*}
And for the part that act on the domains, using the fact that $R_h^{-1} = R'^{-1}R$:
\begin{align*}
      (t_hr_h)^{-1}(tr)^{-1}x_0
      &= R'^{-1}x_0 - R'^{-1}t-R'^{-1}Rt_h \\
      &= R'^{-1}x_0 - R'^{-1}(t+Rt_h) \\
      &= R'^{-1}x_0 - R'^{-1}t' \\
      &= (t'r')^{-1}x_0
\end{align*}

  as well as:
  \begin{align*}
       (t_hr_h)^{-1}(tr)^{-1}(t_1r_1)^{-1} x_0 &=(t'r')^{-1}(t_1r_1)^{-1} x_0
  \end{align*}

Thus, we can continue the computation of the TC:
\begin{align*}
    T(f_2)_{g_1g_2}
        &= \int_{\mathbb{Z}^2} \int_H
            \left(\rho(r)\rho(r_h)
            f((t_hr_h)^{-1}(tr)^{-1}x_0)\right)^\dagger(x_0) \\
                &\qquad\qquad\qquad.\left(
                \rho(r_1)\rho(r)\rho(r_h)
                f((t_hr_h)^{-1}(tr)^{-1}(t_1r_1)^{-1}x_0)\right)\\
                &\qquad\qquad\qquad.\left(
                \rho(r_2)\rho(r)\rho(r_h)
                f((t_hr_h)^{-1}(tr)^{-1}(t_2r_2)^{-1}x_0)\right) dtdr\\
         &= \int_{\mathbb{Z}^2} \int_H
         \left(\rho(r')
            f((t'r')^{-1}x_0)\right)^\dagger(x_0) \\
                &\qquad\qquad\qquad.\left(
                \rho(r_1)\rho(r')
                f((t'r')^{-1}(t_1r_1)^{-1}x_0)\right)\\
                &\qquad\qquad\qquad.\left(
                \rho(r_2)\rho(r')
                f((t'r')^{-1}(t_2r_2)^{-1}x_0)\right) dt'dr'\\
        &=T(f)_{g_1g_2}
\end{align*}
\end{proof}

\subsection{Bispectrum: Definition}



\begin{proof}
We have a group $G = \mathbb{Z}_2 \times H$. We use the convention $g = (t, r) = tr$ for a group element decomposing into a translation component and a rotation (later, of any compact group) one.

By definition, the bispectrum is the Fourier transform of the triple correlation. We have defined the Fourier transform associated with the induced representation's action.

We compute this Fourier transform for the product group $G \times G$ and show that we obtain the expression given above. We recall that $T(f)_{g_1, g_2} \in \mathbb{C}^{K}$. The terms $T_{g}[f]^\dagger T_{g_1g}[f]$ define an inner product and need to stick together.


We define it for any action of $g$ on a vectorial signal, where $g$ can touch both the fiber and the domain. The action is not restricted to the induced representation, thus we use $g$ and not $rt$. A remaining question is whether the induced representation brings something more? Or why does it seem to be important to have the induced representation in the TC but not for the BiSp? It is probably enough as long as the action is homogeneous? Or maybe, in order for it to be an action on fiber+domain it has to be some type of induced representation. Otherwise it does not make sense?
Associativity of mixed tensor-product/matrix multiplication operation
%https://math.stackexchange.com/questions/3029078/associativity-of-mixed-matrix-product-and-tensor-product

The following assumes that: $v \otimes (AB) = (v\otimes A). B$ for $v\in\mathbb{C}^K$ and $A, B \in \mathbb{C}^{D \times D}$

    \begin{align*}
    \beta(f)_{\rho_1, \rho_2}
    &=  \int_{g_1 \in G} \int_{g_2 \in G}
        T(f)_{g_1, g_2} \otimes
            \left[
                \rho_1(g_1)^\dagger \otimes \rho_2(g_2)^\dagger
            \right]
            dg_1 dg_2 \\
    &=
    \int_{g_1 \in G} \int_{g_2 \in G}
        \int_{g \in G}
                \left[ T_{g}[f]^\dagger (x_0) T_{g_1g}[f](x_0) T_{g_2g}[f](x_0)
                    \right]dg
            \otimes
              \left[
                \rho_1(g_1)^\dagger \otimes \rho_2(g_2)^\dagger
              \right]
            dg_1 dg_2 \\
    &=\int_{g \in G}
      \int_{g_1 \in G} \int_{g_2 \in G}
                  \left[T_{g}[f]^\dagger (x_0)T_{g_1g}[f](x_0) T_{g_2g}[f](x_0)
                    \right]
                \otimes
              \left[
                \rho_1(g_1)^\dagger \otimes \rho_2(g_2)^\dagger
              \right]
            dg_1 dg_2dg\\
    &=\int_{g \in G}
      \int_{g_1 \in G} \int_{g_2 \in G}
               T_{g}[f]^\dagger (x_0) T_{g_1}[f](x_0) T_{g_2}[f](x_0)
                \otimes
              \left[
                \rho_1(g_1g^{-1})^\dagger \otimes \rho_2(g_2g^{-1})^\dagger
              \right]
            dg_1 dg_2 dg\\
    &=\int_{g \in G}
      \int_{g_1 \in G} \int_{g_2 \in G}
            [...]
                \otimes
              \left[
                \rho_1(g_1)^\dagger \rho_1(g^{-1})^\dagger \otimes \rho_2(g_2)^\dagger \rho_2(g^{-1})^\dagger
              \right]
            dg_1 dg_2dg\\
        &=\int_{g \in G}
      \int_{g_1 \in G} \int_{g_2 \in G}
            [...]
                \otimes
              \left[
                \rho_1(g_1)^\dagger \rho_1(g) \otimes \rho_2(g_2)^\dagger \rho_2(g)
              \right]
            dg_1 dg_2dg\\
        &=\int_{g \in G}
      \int_{g_1 \in G} \int_{g_2 \in G}
            [...]
                \otimes \left(
              \left[
                \rho_1(g_1)^\dagger \otimes \rho_2(g_2)^\dagger \right] . \left[\rho_1(g) \otimes  \rho_2(g)
              \right] \right)
            dg_1 dg_2dg\\
            &=\int_{g \in G}
      \int_{g_1 \in G} \int_{g_2 \in G}
       \left(T_{g}[f]^\dagger (x_0) T_{g_1}[f](x_0) T_{g_2}[f](x_0)
                \otimes
              \left[
                \rho_1(g_1)^\dagger \otimes \rho_2(g_2)^\dagger \right]\right)  .
            dg_1 dg_2\left[\rho_1(g) \otimes  \rho_2(g)
              \right] dg\\
        &=\int_{g \in G}
        T_{g}[f]^\dagger (x_0) \left(
      \int_{g_1 \in G} \int_{g_2 \in G}
       T_{g_1}[f](x_0) T_{g_2}[f](x_0)
                \otimes
              \left[
                \rho_1(g_1)^\dagger \otimes \rho_2(g_2)^\dagger \right]\right)  .
            dg_1 dg_2\left[\rho_1(g) \otimes  \rho_2(g)
              \right] dg\\
            &=\int_{g \in G}
            T_{g}[f]^\dagger (x_0)
      \left[\mathcal{F}(f)_{\rho_1} \otimes \mathcal{F}(f)_{\rho_2}\right] \left[\rho_1(g) \otimes  \rho_2(g)
              \right] dg\\
            &=\left[\mathcal{F}(f)_{\rho_1} \otimes \mathcal{F}(f)_{\rho_2}\right]
            \int_{g \in G}
            T_{g}[f]^\dagger (x_0)
       \left[\rho_1(g) \otimes  \rho_2(g)
              \right] dg\\
            &=\left[\mathcal{F}(f)_{\rho_1} \otimes \mathcal{F}(f)_{\rho_2}\right]
            \int_{g \in G}
            T_{g}[f]^\dagger (x_0) \otimes
       \left[
        C_{\rho_1, \rho_2}
        \bigoplus_{\rho \in \rho_1 \otimes \rho_2}
                \rho(g)
        C_{\rho_1, \rho_2} ^\dagger
        \right]
              dg\\
            &=\left[\mathcal{F}(f)_{\rho_1} \otimes \mathcal{F}(f)_{\rho_2}\right]
            \int_{g \in G}
            \bigoplus_{\rho \in \rho_1 \otimes \rho_2}
            T_{g}[f]^\dagger (x_0) \otimes
       \left[
        C_{\rho_1, \rho_2}
                \rho(g)
        C_{\rho_1, \rho_2} ^\dagger
        \right]
              dg\\
        &=\left[\mathcal{F}(f)_{\rho_1} \otimes \mathcal{F}(f)_{\rho_2}\right]
        C_{\rho_1, \rho_2}
            \int_{\mathbb{Z}^2}
            \int_H
            \bigoplus_{\rho \in \rho_1 \otimes \rho_2}
            T_{g}[f]^\dagger (x_0) \otimes
       \left[
                \rho(g)
        \right]
              dg C_{\rho_1, \rho_2} ^\dagger \\
        &=\left[\mathcal{F}(f)_{\rho_1} \otimes \mathcal{F}(f)_{\rho_2}\right]
        C_{\rho_1, \rho_2}
            \bigoplus_{\rho \in \rho_1 \otimes \rho_2}
           \mathcal{F}(f)_\rho C_{\rho_1, \rho_2} ^\dagger \\
     \end{align*}
\end{proof}

TO CHECK

\section{Physics-Invariance}
Consider $f$ a solution to a PDE, defined on $\mathbb{R}^2 \times \mathbb{R}$, i.e., on a space-time domain, and with values in some co-domain.

Consider the group $G$ of its symmetries, which are actions on $f$. Most of these actions are: linear actions acting only on the domain of $f$, which is likely homogeneous for the action, thus this fits to one of the cases listed here. Other actions are linear actions, that look like steerable group actions, thus this fits into one of the cases listed here. For the remaining non-linear actions, we can define a Fourier transform, and bispectrum, for this case?

For example, we can think of building a transformer whose attention coef are invariant thanks to the bispectrum. If we do this, does this mean that we should work in the fourier domain, for the  whole network? What does it mean to solve a PDE in Fourier domain?

\section{Scratch}

\subsection{Triple Correlation: Definition (Continuous case/Comp. Cost: Try 1)}

We try to generalize the triple correlation for the case where the subgroup $H$ is continuous, i.e. for when H is a Lie group. We use H = SO(n) as our running example.

The main ingredient for the steerable-TC on is the action of a group. Thus, the trick is to turn this action into the action of the associated Lie algebra, and use the generators of the Lie algebra. The Lie algebra of SO(n) admits three generators that we write $A_1, A_2, A_3$ (in matrix form) or $a_1, a_2, a_3$ (in abstract form).

We show that it is enough to define the TC only on these three generators, to define the whole TC.
\begin{equation}
    T(f)_{a^it_1a^jt_2}
        = \int_{t \in \mathbb{Z}^2}
            \int_{a \in \mathfrak{h}}
                T_{at}[f]^\dagger (x) T_{a^it_1at}[f](x) T_{a^jt_2at}[f](x)
                    dadt,
\end{equation}
for $i, j \in \{1, 2, 3\}$ and any $t_1, t_2 \in \mathbb{Z}^2$. Also, we have: $g = \exp(a)t$, $g_1 = \exp(a^i)t_1$ and $g_2 = \exp(a^j) t_2$ and we define:
\begin{equation}
    T_{at}[f](x) = \rho(\exp(a))f((t\exp(a))^{-1}x_0).
\end{equation}

\begin{lemma}[Group \& algebra representations]
For a lie algebra reps $t(a)\in M_{d\times d}$ of $a \in\mathfrak{h}$ where $r(a) = \exp(a)$ is the corresponding group element:
\begin{equation}
    T_{r(a)} = T_{\exp(a)} = \exp t(a) \in GL_{d \times d}.
\end{equation}
For example, for the representation $\phi$ that appears in the next lemma:
\begin{equation}
    \phi\left(\exp\left(-a\right)\right) = \exp\left(-\sigma(a) \right),
\end{equation}
where $\sigma$ is the lie algebra representation corresponding to the group representation $\phi$.
\end{lemma}

\begin{lemma}[Action of an inverse element on $x$]
    For $a \in \mathfrak{h}$ where $a = \sum_{l=1}^3 w_l a^l$ and $r = \exp(a)$, we have:
    \begin{align*}
        \left(tr\right)^{-1}x
        &= \left(
            r^{-1}, -R^{-1} t\right)x \\
        &= \left(
            \exp\left(a\right)^{-1}, -\phi\left(\exp\left(a\right)\right)^{-1} t\right)x \\
        &= \left(
            \exp\left(-a\right), -\phi\left(\exp\left(-a\right)\right) t\right)x \\
        &= \left(
            \exp\left(-a\right), -\exp\left(-\sigma(a) \right) t\right)x \\
        &= \exp\left(-a\right)x - \exp\left(-\sigma(a) \right) t
    \end{align*}
\end{lemma}

In the following lemma, we use convenient notations that will help for the last derivation of this page.

\begin{lemma}[Action on the signal $f$]
    For $a_2 \in \mathfrak{h}$ where $a_2 = \sum_{l=1}^3 w_l a^l$, and a function signal $F$, we have:
    \begin{align*}
        T_{a_2t_2}[F](x)
         &= \rho(\exp(a_2))F((t_2\exp(a_2))^{-1}x)\\
         &= \exp(\tau(a_2)) F\left[\exp\left(-a_2\right)x - \exp\left(-\sigma(a_2) \right) t_2\right] \\
         &= \exp\left(\tau\left(\sum_{l=1}^3 w_l a^l\right)\right)
            F\left[
                \exp\left(-\sum_{l=1}^3 w_l a^l\right)x
                - \exp\left(-\sigma(\sum_{l=1}^3 w_l a^l) \right) t_2
            \right] \\
         &= \exp\left( -\sum_{l=1} w_l \sigma(a^l)\right)
                F\left[
                    \exp\left(-\sum_{l=1}^3 w_l a^l\right)x
                    - \exp\left(-\sigma(\sum_{l=1}^3 w_l a^l) \right) t_2
                    \right]
    \end{align*}

For the next computations, we assume that the group $H$ is commutative, so that $\exp(A+B) = \exp(A)\exp(B).$
\begin{align*}
    T_{a_2t_2}[F](x)
    &= \exp\left( -\sum_{l} w_l \sigma(a^l)\right)
            F\left[
                \exp\left(-\sum_{l} w_l a^l\right)x
                - \exp\left(-\sigma(\sum_{l} w_l a^l) \right) t_2
                \right]\\
    &= \Pi_{l}\exp\left( -w_l \sigma(a^l)\right)
            F\left[
                \Pi_{l}\exp\left(-w_l a^l\right)x
                - \exp\left(-\sigma(\sum_{l} w_l a^l) \right) t_2
                \right]
    \end{align*}
This seems too complicated. We start again. We note that:
\begin{align*}
    T_{a_2t_2}[F](x)
     &= T_{\exp(a_2)t_2}[F](x) \\
     &= T_{\exp(\sum_{l=1}^3 w_l a^l)t_2}[F](x) \\
     &= T_{\Pi_{l=1}^3\exp(w_l a^l)t_2}[F](x) \\
     &= T_{\Pi_{l=1}^3\left(\exp(a^l)\right)^{w_l}t_2}[F](x) \\
     &= \Pi_{l=1}^3\left(T_{\exp(a^l)}\right)^{w_l}T_{t_2}[F](x) \\
\end{align*}
\end{lemma}

Now, we show that any $T(f)_{a_1t_1a_2t_2}$ can be expressed as a function of $T(f)_{a_it^1a^jt_2}$. By symmetry in $a_1, a_2$, we only prove it for $T(f)_{a^it_1a_2t_2}$.

For $a_2 \in \mathfrak{h}$ where $a_2 = \sum_{l=1}^3 w_l a^l$, we have:
\begin{align*}
    T(f)_{a^it_1a_2t_2}
    &= \int_{t \in \mathbb{Z}^2}
        \int_{a \in \mathfrak{h}}
            T_{at}[f]^\dagger (x) T_{a^it_1at}[f](x) T_{a_2t_2at}[f](x)
                dadt \\
    &= \int_{t \in \mathbb{Z}^2}
        \int_{a \in \mathfrak{h}}
            T_{at}[f]^\dagger (x) T_{a^it_1at}[f](x) T_{a_2t_2}[T_{at}(f)](x)
                dadt \\
    &=\int_{t \in \mathbb{Z}^2}
        \int_{a \in \mathfrak{h}}
            F^\dagger (x) T_{a^it_1}[F](x) T_{a_2t_2}[F](x)
                dadt \quad \text{notation: $F = T_{at}(f)$}\\
  &=\int_{t \in \mathbb{Z}^2}
        \int_{a \in \mathfrak{h}}
            F^\dagger (x) T_{a^it_1}[F](x) \Pi_{l=1}^3\left(T_{\exp(a^l)}\right)^{w_l}T_{t_2}[F](x)
                dadt
\end{align*}
Here, we are blocked because we cannot unpack the product $\Pi_l$.

\subsection{Triple Correlation: Definition (Continuous case/Comp. Cost: Try 2)}

We consider the TC as an operator:
\begin{align*}
    TC: \mathcal{F} \times \mathfrak{g}^2 &\mapsto \mathbb{R}^K \\
        f, (t_1a_1, t_2 a_2) &\mapsto T(f)_{t_1a_1t_2a_2}
\end{align*}
This is not a linear operator, thus it cannot be represented by a big tensor. This is the main difference with the case of the convolution: the convolution being a linear operator, we can represent it as a big tensor that contains the whole filter bank. The nonlinearity in $f$ comes from the triple product. The nonlinearity in $a_2$ comes from the fact that we were not able to prove it in Try 1.

Let's try again though. With the discrete group case for now. In the convolution case, all the rotated filters are included in the filter bank. If we think in the same way for the TC, this means that we should see the TC as:
\begin{align*}
    TC: \mathcal{F}&\mapsto \mathbb{R}^{K \times \|G\|^2} \\
        f  &\mapsto \left[T(f)_{t_1a_1t_2a_2}, ..., ... \right],
\end{align*}
where on the output we have all the values of the TC for all the (non redundante) pairs of group elements.

We have $\mathcal{F} = \{f:\mathbb{Z}^2 \mapsto \mathbb{R}^K \} = \mathbb{R}^{S^2.K} $ where $S$ is the side of the $\mathbb{Z}^2$ grid that supports the image $f$.

As an operator on $f$, what is the polynomial order of the TC?

\begin{align*}
    TC(\mu f+\lambda g)
    &= \int_{t \in \mathbb{Z}^2}
        \int_{a \in \mathfrak{h}}
            T_{at}[\mu f+\lambda g]^\dagger (x)
            T_{a_1t_1at}[\mu f+\lambda g](x)
            T_{a_2t_2at}[\mu f+\lambda g](x)
                dadt \\
    &= \int_{t \in \mathbb{Z}^2}
        \int_{a \in \mathfrak{h}}
            \left(\mu T_{at}[f] + \lambda T_{at}[g]\right)^\dagger (x)
            \left(\mu T_{a_1t_1at}[f] + \lambda T_{a_1t_1at}[g]\right) (x)
            \left(\mu T_{a_2t_2at}[f] + \lambda T_{a_2t_2at}[g]\right) (x)
                dadt \\
    &= ...\text{not equal to $ \mu TC(f) +\lambda TC(g)$, nor with its higher orders, bcs we have crossed terms with $f, g$}
\end{align*}

Thus, TC is just an invariant function on $f$. Since it is not linear, we cannot represent it as a matrix and write the intertwiner equation of the type $T \pi(g) = \rho(g) T$ that was useful to find a basis for the filters in the G-conv of the steerable case.

\subsection{Triple Correlation: Definition (Continuous case/Comp. Cost: Try 3)}

The last idea would be to think of the triple correlation as an implicit representation. For example, we could implement it as a continuous function (a neural network) takes a pair $(g_1, g_2)$ as input and output the $T(f)_{g_1g_2}$. The NN could implement the ansatz of the triple correlation, e.g. have a triple product structure that is able to evaluate the TC each time $f, g_1, g_2$ is given. The weights of this neural network would contain everything we need about the TC, without having to evaluate it at each pair of group elements. However, it is possible that we only shift the problem is that the way of implementing the TC as an implicit function is going to be even more computationally expensive that computing all of the values of the TC.

\subsection{More comments}

The computational cost mostly comes from the translation part, i.e. from the size of the grid of the images. We note that, for steerable CNNs, only the abstraction related to the group H in the semi direct product $G = H \times Z^2$ is used in the filter bank coefficients. For the translation part, there is no computational gain: the filter bank is still moved at each position during the convolution. This could be an argument in favor of not trying to use the steerable strategy for complexity gain.

\subsection{More comments: continuous case}

Since the obstruction to using the steerable strategy is also the fact that the TC is non linear, we could try to go back to a situation where we have a linear operator. For example, we could consider the bispectrum which in the discrete setting is of the form
\begin{equation}
    \beta(x) = Ux \otimes Vx \otimes Wx
\end{equation}
where $U, V, W$ are different realization of the discrete Fourier transform.

Considering the case of the continuous fourier transform instead, we know that the Fourier transform is an equivariant linear operator:
\begin{equation}
    W\phi(g) = \rho(g) W,
\end{equation}
when $W$ is the discrete fourier matrix. This equation should hold for the case where the image is discrete (discrete in translation) but H is continuous. The continuous fourier transform in H could be decomposed into a discrete set of coefficients on a equivariant basis of intertwiners. This could make the (fixed, not learnable) bispectrum neural network applicable to continuous groups.

We need the fourier transform for the steerable case, though.


\bibliographystyle{plos2015}
\bibliography{references.bib}

\end{document}
